\documentclass[11pt]{article}
\usepackage[T1]{fontenc}
\usepackage[utf8]{inputenc}
\usepackage{lmodern,amsmath,amssymb,amsthm,mathtools}
\usepackage[margin=1in,headheight=14pt]{geometry}
\usepackage{microtype,booktabs,longtable,array,enumitem,listings,xcolor,fancyhdr,xurl}
\usepackage[colorlinks=true,linkcolor=blue!45!black,urlcolor=blue!45!black,citecolor=blue!45!black]{hyperref}
\newtheorem{theorem}{Theorem}[section]
\newtheorem{lemma}[theorem]{Lemma}
\newtheorem{proposition}[theorem]{Proposition}
\newtheorem{corollary}[theorem]{Corollary}
\theoremstyle{definition}
\newtheorem{definition}[theorem]{Definition}
\newcommand{\ipra}{\mathsf{iPRA}}
\newcommand{\C}{\mathcal C_5}
\newcommand{\NZ}{\mathcal N_0}
\newcommand{\nf}{\operatorname{nf}}
\newcommand{\fnf}{\operatorname{fnf}}
\newcommand{\si}{\mathcal I}
\newcommand{\Eq}{\mathsf{Eq}}
\newcommand{\UE}{\mathsf{Eval}}
\newcommand{\RF}{\mathsf{Refl}}
\newcommand{\Q}{\mathsf Q}
\newcommand{\D}{\mathsf D}
\newcommand{\syn}{\equiv}
\newcommand{\FV}{\operatorname{FV}}
\newcommand{\file}[1]{\expandafter\nolinkurl\expandafter{\detokenize{#1}}}
\newcommand{\code}[1]{\texttt{#1}}
\setlist{nosep}
\setlength{\emergencystretch}{2em}
\pagestyle{fancy}
\fancyhf{}
\fancyhead[L]{\small Five-rule iPRA with one evaluation schema}
\fancyhead[R]{\small Research report}
\fancyfoot[C]{\thepage}
\lstset{basicstyle=\ttfamily\footnotesize,breaklines=true,columns=fullflexible,
  keepspaces=true,showstringspaces=false,frame=single,rulecolor=\color{black!15},
  xleftmargin=3pt,xrightmargin=3pt,aboveskip=8pt,belowskip=8pt}
\hypersetup{pdftitle={Five-rule iPRA with a single evaluation schema},pdfauthor={Research report prepared for Vladimir Reshetnikov}}
\title{\textbf{Five-rule iPRA with a single evaluation schema}\\[5pt]
\large Clause contexts, zero normalization, and a smaller complete verifier}
\author{Prepared for Vladimir Reshetnikov}
\date{November 11, 2026}
\begin{document}
\maketitle
\begin{abstract}
We give a complete presentation of full first-order intuitionistic primitive
recursive arithmetic with two formula constructors, equality and a paired
binder, and five inference patterns: reflexivity, one uniform evaluation
schema, clause use, paired elimination, and paired introduction. Zero
iteration is contracted during term construction and substitution; successor
iteration is retained as syntax. The evaluation schema performs both
computation and all equality reasoning, with no separate equality-only
context parameter and no case split. The presentation is bi-interpretable and equiconsistent
with iPRA, with whole numerical domains, actual equality, and numerical
composite graph $x=y$.

The construction combines two existing reductions in the current repository:
single successor evaluation on zero-normal terms and clause contexts that
remove first-class conditional closure. A shorter two-evaluation congruence
proof improves the accompanying proof compiler. The complete verifier is
measured against controls with the same input discipline, including parsing,
validation, normalization, substitution, and inference. The reduction is
strict in all three source-size coordinates against the matched raw control.
It does not establish an absolute minimum or beat the repository's separate
516-byte compact codec. The new equivalence is a written proof, with
executable proof translations and independent implementation checks; it is
not a new Lean or Rocq formalization.
\end{abstract}
\clearpage
\tableofcontents
\newpage

\section{Result and comparison boundary}

The recommended presentation is $\C$. Its five primitive inference patterns
are
\[
\boxed{\RF,\quad\UE,\quad\mathsf{Use},\quad\mathsf{Elim},\quad\mathsf{Intro}.}
\]
The only formula constructors are $s=t$ and $\D_{x,y}B$. There is one
numerical sort. Its terms are variables, $0$, successor, and scalar iteration,
subject only to the zero normalization specified below. The single
evaluation schema is
\[
\frac{\Gamma\vdash\si(n,a;u.h)=t\qquad
      \Gamma\vdash I(Sn,a;u.h)=r}
     {\Gamma\vdash r=h\langle t/u\rangle}.
\tag{Eval}\label{eq:main-eval}
\]
Here $\si(0,a;u.h)=a$ by syntax, and angle brackets denote capture-avoiding
substitution followed by zero normalization. Nothing else is evaluated by
the term representation. In particular, $I(S0,a;u.h)$ remains an iterator
node. This surviving one-step iterator supplies the complete equality
interface.

\begin{theorem}[Main result]\label{thm:main}
The ordinary realization of $\C$ is bi-interpretable and equiconsistent
with full first-order $\ipra$. Both interpretations use the whole numerical
domain and actual equality. Both numerical composites are internally
isomorphic to the identity by the graph $x=y$. The formula translations in
both directions have provable round trips. The result holds for arbitrary
finite retained contexts of ordinary formulas, not only closed theorems.
\end{theorem}

\subsection{What is being improved}

The comparison is pinned to ProveIt commit
\file{1eb151b8f40800aec569361e827d934cb77bd3e1}. Its index describes
52 filed packages, and its incoming directory contains four more archives.
The incoming packages are relevant: ignoring them would overstate novelty.
The closest proposals are the following.

\begin{center}
\begin{tabular}{@{}p{.25\linewidth}p{.19\linewidth}p{.44\linewidth}@{}}
\toprule
Presentation & Formula syntax & Arithmetic/equality interface\\
\midrule
Raw clausal CE5 & Eq/$D$ & Refl and three-case Eval; ordinary substitution\\
Zero-normal $\mathcal Z_1$ & Eq/$Q$/$D$ & Refl and one successor Eval; six patterns\\
Canonical clausal & Eq/$D$ & Successor-seeded recursor, computation and supplied transport template\\
Present $\C$ & Eq/$D$ & Refl and one successor Eval; five patterns\\
\bottomrule
\end{tabular}
\end{center}

The mathematical component reductions are inherited and are credited to
their repository packages~\cite{rawclausal,iteration,clausal,baseline}.
The contribution here is their exact five-pattern combination, the explicit
comparison proof on normal syntax, the shorter congruence compiler, and a
complete controlled implementation. No previously known component is
presented as a new axiom-deletion theorem. The complete single-evaluation
clause combination was not present in the inspected current specifications;
this is a comparison with that snapshot, not a literature-priority claim.

The improvement removes actual kernel work: the raw evaluator's root,
zero-iteration, and successor-iteration alternatives become one schema.
It also removes the native $Q$ formula constructor and its dedicated
hypothesis/elimination interface. The normalization traversal and clause
representation are explicitly counted. Five patterns alone are not a
minimality measure: combining unrelated rules under one name would prove
nothing about simplicity.

\subsection{Measured complete checkers}

The selected standalone program has 2,283 normalized UTF-8 bytes,
741 significant tokens, and 826 AST nodes. Its matched raw control has
2,416 bytes, 744 tokens, and 842 nodes. Thus the byte reduction is
133 bytes (5.50\%); the token and AST reductions are 3 and 16.
\begin{center}
\begin{tabular}{@{}lrrrr@{}}
\toprule
Complete goal-directed program & Bytes & Tokens & AST nodes & Gzip-9\\
\midrule
Raw CE5 control & 2,416 & 744 & 842 & 855\\
Zero-normal control with root case & 2,426 & 765 & 865 & 862\\
\textbf{Single-Eval clause checker} & \textbf{2,283} & \textbf{741} & \textbf{826} & \textbf{818}\\
\bottomrule
\end{tabular}
\end{center}
All entries count the entire acceptance module after removing its module
docstring, applying CPython 3.12.14 \file{ast.unparse}, and appending one LF.
Significant tokens exclude whitespace, indentation, comments, and the end
marker. AST nodes are counted by \file{ast.walk}. Gzip uses level 9 and
zero timestamp and is informational, not an optimization objective.

The readable selected source is 2,717 bytes, or 1,036 bytes with gzip-9.
The readable raw control is 2,698 bytes, or 984 with gzip-9; thus the
normalized-source comparison is not a claim that the original formatted
file is shorter. The auxiliary forward implementation is 3,022 normalized
bytes, 867 tokens, and 1,008 AST nodes, against 3,046, 853, and 993 for
its forward control. It improves bytes only; it is not the selected
coordinatewise result.

The main comparison uses the same goal-directed protocol, clause contexts,
variable representation, validated reflexive leaves, and public output
contract. The successor-only rule naturally supplies the parameters of its
two premises, rather than testing several possible shapes after both proofs
have been checked. That certificate representation is part of the specified
presentation and is included in the measured program.

The historical compact codec has a different representation and a highly
compressed proof protocol. The present complete checker is larger than its
516-byte normalized source. Thus this report establishes a structural and
matched implementation improvement; it does not assert a new global lower
envelope. There is also no proof that a globally simplest axiomatization
exists for an unspecified implementation language and coding convention.

\section{The entire formal theory}
\label{sec:specification}

\subsection{Normal terms and substitution}

Start with the raw syntax
\[
 t::=x\mid0\mid S(t)\mid I(n,a;u.h).
\]
Only $u$ in the update body $h$ is bound. The counter and initializer
are outside that binding. Updates may have any fixed numerical parameters.
There is no variable for a program or function description and no universal
evaluator. Define $\nf$ recursively on this finite syntax:
\[
\begin{aligned}
\nf(x)&=x,&\nf(0)&=0,&\nf(St)&=S\nf(t),\\
\nf(I(n,a;u.h))&=
\begin{cases}
\nf(a),&\nf(n)\syn0,\\
I(\nf(n),\nf(a);u.\nf(h)),&\text{otherwise}.
\end{cases}
\end{aligned}
\]
Bound names are first renamed when needed. A term of $\C$ is a fixed
point of $\nf$: every iterator node has a counter that is not literally
$0$. On normal inputs write
\[
\si(n,a;u.h)=
\begin{cases}a,&n\syn0,\\ I(n,a;u.h),&n\not\syn0.\end{cases}
\]
The normalization is syntactic. A counter provably equal to zero is not
contracted unless it is the literal zero after the recursive normalization.
Normal substitution is
\[
       t\langle s/x\rangle:=\nf(t[s/x]).
\]
Formulas and clauses use the corresponding operation on every contained
term. The full successor computation is deliberately not made definitional.

\begin{lemma}[Zero normalization]\label{lem:normalization}
Zero normalization terminates, preserves scope, and is idempotent. It
commutes with capture-avoiding substitution in the form
\[
\nf(t[s/x])=\nf\bigl(\nf(t)[\nf(s)/x]\bigr).
\]
Normal substitutions obey the usual composition identities. For each raw
term, the ordinary iterator theory proves $t=\nf(t)$.
\end{lemma}
\begin{proof}
The definition recurses only on proper subterms and contracts an entire
zero-counter iterator to its already normalized initializer. Induction on
syntax proves the first three assertions and the displayed identity. The
iterator case distinguishes whether its normalized substituted counter is
literal zero; both sides then make the same choice. The ordinary substitution
composition law followed by this identity gives the normal version.

For provable equality, use the zero computation equation and congruence
on the counter and seed. For changes inside the update body, equation
induction proves extensionality of a fixed iterator: equal seeds and
pointwise equal updates give equal iterates. This handles the binder body
without treating it as an unmentioned ordinary function argument.
The repository's normalization development proves this comparison formally
for the reference kernel~\cite{normal}.
\end{proof}

\subsection{Formulas, retained clauses, and falsity}

The complete formula grammar is
\[
        A::=(s=t)\mid\D_{x,y}A.
\]
The paired constructor binds two distinct numerical variables. A context
is a finite reusable list of clause occurrences
\[
      [A;y.B(y)],\qquad y\notin\FV(A).
\]
The clause binds $y$ only in its consequent. It is an object of the context,
not a third formula constructor. An ordinary assumption $F$ abbreviates
\[
       \operatorname{ord}(F):=[H;v.F],
       \qquad H:=(0=0),\quad v\notin\FV(F).
\]
Accordingly, $\Gamma,F$ in a rule below means
$\Gamma,\operatorname{ord}(F)$. Clause occurrences are distinguished,
even when two occurrences have identical text.
The designated contradiction is $0=S0$.

\clearpage
\subsection{All five primitive inference patterns}

All numerical data and formulas must be well scoped and normal. Every
instance uses the normal substitution just defined.
\[
\frac{}{\Gamma\vdash s=s}.\tag{Refl}
\]
\[
\frac{\Gamma\vdash\si(n,a;u.h)=t\qquad
      \Gamma\vdash I(Sn,a;u.h)=r}
     {\Gamma\vdash r=h\langle t/u\rangle}.\tag{Eval}
\]
\[
\frac{\Gamma\vdash A}{\Gamma\vdash B\langle t/y\rangle}
\quad\text{when }[A;y.B]\text{ is retained in }\Gamma.\tag{Use}
\]
\[
\frac{\Gamma\vdash\D_{x,y}B(x,y)\qquad
      \Gamma,[B(b,0);y.B(b,y)]\vdash F}
     {\Gamma\vdash F}.\tag{Elim}
\]
\[
\frac{\Gamma,B(t,\eta)\vdash B(t,Sz)}
     {\Gamma\vdash\D_{x,y}B(x,y)},\qquad
\eta=\begin{cases}
z,& B\text{ is a literal equation and }t\syn0,\\
0,&\text{otherwise}.
\end{cases}\tag{Intro}
\]
In Elim, $b$ is fresh for $\Gamma$, $F$, and all passive parameters of
$B$. In Intro, $z$ is fresh for $\Gamma$, $t$, and all passive parameters
of $B$. The displayed local clause or ordinary assumption is discharged.
The update variable $u$ in Eval binds only $h$. No freshness condition
excludes passive parameters from that update.

There are no further logical, structural, equality, successor, computation,
or induction axioms or rules. In particular, there is no primitive
hypothesis rule, no $Q$ formula constructor, no arbitrary formula induction,
and no separate equality-transport rule. Eval does supply the iterator
update $h$, whose occurrences are checked in its computation premises;
there is no additional equality-only context parameter. The mode test in
Intro is part of the specification, not an unstated inference choice.

\section{Recovering equality and the missing case introduction}

\begin{lemma}[Elementary structural operations]\label{lem:structure}
Ordinary hypothesis use, weakening, injective fresh renaming, and cut for
ordinary formula assumptions are admissible in $\C$.
\end{lemma}
\begin{proof}
Refl proves $H$, and Use on $[H;v.F]$ proves $F$. For weakening and
renaming, induct on a derivation, first freshening eigenvariables away from
new retained data. Injective renaming preserves normal terms, equation roots,
and literal zero, so it preserves Intro's mode.

For cut, replace a Use of the distinguished occurrence $[H;v.F]$ by
the supplied proof of $F$, weakened into its local context. At other rules
transform the premises recursively. The extra proof of $H$ at that Use
can be discarded. No arbitrary numerical substitution into a proof is
assumed in this argument.
\end{proof}

\begin{lemma}[Identity frame and equality composition]\label{lem:root}
From $p:s=t$ and $q:s=r$, $\C$ derives $r=t$ using two Eval nodes
and one new reflexivity leaf, with each supplied proof used once.
\end{lemma}
\begin{proof}
Let $K(s)=I(S0,s;u.u)$. Its predecessor iterate is definitionally
$\si(0,s;u.u)=s$. Thus
\[
\operatorname{Root}(p,q):=
\UE\bigl(p,\UE(q,\RF(K(s)))\bigr):r=t.
\]
The inner Eval proves $K(s)=r$, and the outer one proves $r=t$.
No equality law or logical principle is used beforehand.
\end{proof}

Symmetry is $\operatorname{Root}(\RF(s),p):t=s$ for $p:s=t$.
Transitivity of $p:s=t$ and $q:t=v$ is
$\operatorname{Root}(q,\operatorname{Sym}(p)):s=v$.
Successor computation is a single Eval on two reflexivities:
\[
 I(Sn,a;u.h)=h\langle\si(n,a;u.h)/u\rangle.
\]
Zero computation is already an identity of the normal syntax.

\begin{lemma}[Congruence in two evaluations]\label{lem:short-congruence}
For every normal numerical template $h(u)$, from $p:s=t$ one can
derive $h\langle s/u\rangle=h\langle t/u\rangle$ by two Eval nodes
and two reflexivity leaves. The input proof occurs once.
\end{lemma}
\begin{proof}
Put $J=I(S0,s;u.h)$. First construct
\[
 e_s:=\UE(\RF(s),\RF(J)):J=h\langle s/u\rangle.
\]
The left side of this equation is still the successor iterator $J$.
Consequently a second application,
\[
          \UE(p,e_s):h\langle s/u\rangle=h\langle t/u\rangle,
\]
is an instance of exactly the same schema. The first premise now uses
the supplied equality in place of the reflexivity of the seed. This
avoids constructing two separately evaluated endpoints and joining them
afterwards.
\end{proof}

\begin{corollary}[Template transport and exact local overhead]
\label{cor:transport}
From $p:s=t$ and $q:h\langle s/u\rangle=r$, derive
$r=h\langle t/u\rangle$ with four Eval nodes and three reflexivity
leaves. Each of $p,q$ occurs once. Replacing $k$ such one-term transport
nodes in a proof whose other rules are retained adds exactly $6k$
inference nodes.
\end{corollary}
\begin{proof}
Apply $\operatorname{Root}$ to the congruence of the preceding lemma
and $q$. The replacement has seven new nodes in place of the old one,
and it does not duplicate either premise tree. This counts inference
nodes, not the serialized sizes of repeated term labels.
\end{proof}

Transport in both sides of an equation follows from two crossed one-term
transports: $a(s)=b(s)$ first gives $b(s)=a(t)$, then $a(t)=b(t)$.
No tree-size bound that hides the two uses of the equality proof is claimed.
The shorter congruence construction is also used by the executable
certificate compiler.

\begin{lemma}[Unrestricted case introduction]\label{lem:cases}
For every matrix $B$ and every normal witness $t$,
\[
\frac{\Gamma,B(t,0)\vdash B(t,Sz)}
     {\Gamma\vdash\D_{x,y}B(x,y)}
\]
is admissible, with the usual freshness condition on $z$.
\end{lemma}
\begin{proof}
Only a literal equation matrix at witness $0$ is excluded by the primitive
case mode. In that situation set
\[
                  Z=I(S0,0;u.0).
\]
This is a valid normal term, is not literally zero, and Eval proves
$Z=0$. Under the local equation $B(Z,0)$, use equation-template
transport to derive $B(0,0)$. Cut in the supplied case proof to obtain
$B(0,Sz)$, then use the reverse equality to derive $B(Z,Sz)$.
Intro at witness $Z$ is in case mode and has the desired conclusion
$D_{x,y}B(x,y)$. Its witnessing term has changed; its conclusion has not.

The raw alias $I(0,0;u.0)$ cannot be used here: it normalizes to $0$.
The one-step alias $Z$ survives because successor computation is not a
normalization rule. This is the relevant interaction between the two
component reductions.
\end{proof}

\begin{lemma}[Transport through all formulas]\label{lem:formula-transport}
For every Eq/$D$ formula template $F(v)$, $s=t$ and $F(s)$ imply
$F(t)$ in the same retained context.
\end{lemma}
\begin{proof}
Use external structural induction on $F$. The equation case was proved
above. For $F(v)=D_{x,y}B(v,x,y)$, eliminate $F(s)$ at fresh $b$,
opening $[B(s,b,0);y.B(s,b,y)]$. Introduce $F(t)$ by the unrestricted
case rule with witness $b$. Convert its local base $B(t,b,0)$ to
$B(s,b,0)$ by the induction hypothesis and symmetry. Use the opened
clause at $Sz$ and convert forward to $B(t,b,Sz)$. Close introduction
and elimination. Every recursive transport concerns the smaller matrix.
\end{proof}

\section{Exact consequence equivalence}

\subsection{Conditional closure is a macro}

For an unused fresh numerical variable $w$, put
\[
                       q_yB:=D_{w,y}B.
\]
This has the ordinary reading $B(0)\to\forall yB(y)$, since the
dummy existential witness ranges over an inhabited numerical domain.

\begin{lemma}[Closure interfaces]\label{lem:qinterfaces}
The following three constructions are admissible:
\begin{enumerate}
\item From $\Gamma,B(0)\vdash B(Sz)$, derive $\Gamma\vdash q_yB$.
\item From $\Gamma\vdash q_yB$ and $\Gamma\vdash B(0)$, derive
      $\Gamma\vdash B(t)$ for every normal $t$.
\item In a context containing $[B(0);y.B(y)]$, derive $q_yB$.
\end{enumerate}
For an equation $E(y)$, one also derives $q_yE$ from
$\Gamma,E(z)\vdash E(Sz)$.
\end{lemma}
\begin{proof}
The first construction is case introduction with an unused witness.
For the second, eliminate the dummy existential and Use the resulting
clause with the supplied base. The opened witness does not occur in $B$.
For the third, use case introduction and specialize the retained clause
at $Sz$ in its branch. For equation progress, use primitive Intro with
dummy witness $0$; it selects the progress assumption $E(z)$.
\end{proof}

\subsection{Comparison with the zero-normal native kernel}

Let $\NZ$ denote the repository's zero-normal kernel, \file{NZPrf}.
It has Eq/$Q$/$D$ formulas and ordinary formula contexts. Its rules
are Hyp, successor computation, crossed one-term transport, $Q$
elimination, $D$ elimination, unrestricted case introduction, and
equation-progress introduction of $Q$. The mathematical comparison is
made with that kernel, rather than assuming that witness-selected Intro
is preserved directly by all numerical substitutions.

The Lean encoding permits raw numerical rule parameters in some places.
Normalize those parameters first. The normalization/substitution identities
preserve the resulting normal premises and conclusion, so the canonical-data
presentation used here has exactly its normal sequents.

Define $T$ on its formulas by retaining equality and $D$, and replacing
every $Q_yB$ by $D_{w,y}T(B)$ with a fresh unused $w$. Let $U$ include
canonical Eq/$D$ formulas into the native formula grammar. Numerical
terms are unchanged by these two maps.

\begin{theorem}[Exact normal-sequent comparison]\label{thm:normal-equiv}
For every finite normal native context $\Gamma$ and normal native formula $A$,
\[
 \Gamma\vdash_{\NZ}A
 \quad\Longrightarrow\quad
 \operatorname{ord}(T\Gamma)\vdash_{\C}T(A).
\]
For normal Eq/$D$ contexts $\Delta$ and normal formulas $F$,
\[
 \operatorname{ord}(\Delta)\vdash_{\C}F
 \quad\Longleftrightarrow\quad
 U\Delta\vdash_{\NZ}UF.
\]
Both maps are effective transformations of finite proof trees.
\end{theorem}
\begin{proof}
Translate a native proof recursively. Hyp is Refl at $0$ followed by
Use. Successor computation and crossed transport are derived above.
The two native introductions are supplied by the case lemma and the
progress clause in Lemma~\ref{lem:qinterfaces}. Native $Q$ elimination
uses its second interface.

For native $D$ elimination, translate the major proof and open the
corresponding clause. The third interface derives $q_yT(B)(b,y)$ in
that clause context, which is the translation of the old $Q$ assumption.
Use ordinary-assumption cut to insert this proof into the translated
branch, then close Elim. This avoids treating a context clause as if it
were already a formula assumption.

Conversely, track the origin of each retained clause occurrence. An
ordinary clause $[H;v.G]$ is represented by native assumption $UG$.
A clause introduced by Elim is represented by $Q_yUB(b,y)$. Use of
the former is Hyp, while Use of the latter is native $Q$ elimination.
Refl is derivable natively. Eval follows by successor computation,
transport along its first equality, and the reverse of its second
equality. Elim becomes native $D$ elimination.

Intro in case mode is the native case rule. In progress mode the matrix
is an equation; native equation progress yields its conditional closure,
which provides the branch for native case introduction at the specified
witness. All freshness conditions are preserved. Initial ordinary
contexts and every locally introduced clause have the stated origin,
so induction on the proof gives the converse for all such sequents.
\end{proof}

\begin{corollary}[Numerical substitution on derivations]
Capture-avoiding normal numerical substitution is admissible for ordinary
retained-formula sequents of $\C$.
\end{corollary}
\begin{proof}
Translate into the stable zero-normal native kernel, apply its admissible
normal substitution, and translate back. The formula translations commute
with normal substitution up to fresh bound renaming. This argument permits
a witness to become literal zero without assuming that Intro's mode is
preserved rule for rule. It supplies the structural operation used in the
subsequent logical macro derivations.
\end{proof}

Arbitrary externally retained clauses have their ordinary reading
$A\to\forall yB(y)$. They can be represented by the native implication
and universal macros; clause use becomes the corresponding two ordinary
eliminations. Thus the optional external clause notation adds no arithmetic
or logical strength.

\begin{theorem}[Formula round trips]\label{thm:roundtrips}
For every normal Eq/$D$ formula $F$, $TUF=F$ literally up to bound
renaming. For every normal native formula $A$,
\[
                   A\dashv\vdash_{\NZ}UT(A).
\]
For raw native formulas, normalize all terms first. The inclusion back into
raw syntax is provably equivalent to the original formula.
\end{theorem}
\begin{proof}
The only nontrivial normal formula step is
$Q_yB\dashv\vdash D_{w,y}B$ for unused $w$. From the left, use
case introduction and closure elimination. From the right, $D$
elimination opens exactly $Q_yB$, which closes the branch by Hyp.
Pointwise interderivability of matrices lifts through $D$ by elimination,
case introduction, backward conversion of its zero assumption, and forward
conversion of its successor goal. It lifts through $Q$ by the same
unused-witness argument. External structural induction gives the result.

For the raw statement, compose with the zero-normalization comparison
\file{nzprf_iff_nprf} and the provable normal-form equivalence
in the pinned development~\cite{normal}. Normalization is part of the
translation; no unproved substitution theorem for witness-selected Intro
is needed.
\end{proof}

\section{Full iPRA and both interpretation composites}
\label{sec:interpretation}

\subsection{The ordinary intermediate theory}

The ordinary scalar theory $T_I$ is intuitionistic first-order logic with
actual equality, $0$, successor, and fixed finite scalar-iterator
descriptions. Each update $h(u,\vec p)$ has defining equations
\[
 I(0,a;u.h)=a,\qquad I(Sn,a;u.h)=h[I(n,a;u.h)/u].
\]
Its other nonlogical principles are coverage,
$\forall x(x=0\lor\exists y\,x=Sy)$, separation
$0=S0\to\bot$, and induction for literal equations with arbitrary
retained assumptions and the standard freshness condition. The native
kernel is a presentation of this ordinary theory, including its logic;
it is not merely an equational fragment~\cite{bridge}.

The readings of our new syntax are
\[
 (s=t)^\circ=(s=t),\qquad
 (D_{x,y}B)^\circ=\exists x\bigl(B^\circ(x,0)\to\forall yB^\circ(x,y)\bigr),
\]
\[
                 [A;y.B]^\circ=A^\circ\to\forall yB^\circ(y).
\]
All five primitive rules are sound under these readings. In particular,
progress Intro is sound because its matrix is literally an equation;
case Intro uses numerical coverage. Using progress for an arbitrary
compound formula would be an unjustified strengthening.

\subsection{Explicit logical macros}

For normal terms let $C(v,a,b)=\si(v,a;u.b)$, with $u$ unused in $b$.
Its zero value is definitionally $a$ and its successor value is provably
$b$. Define a gate recursively by
\[
G_v(s=t):=(C(v,0,s)=C(v,0,t)),\qquad
G_v(D_{x,y}B):=D_{x,y}G_vB,
\]
freshening bound variables away from $v$. Then $G_0A$ is a theorem,
$G_{Sz}A\dashv\vdash A$, and the padding
$D_{u,y}G_yA$ is interderivable with $A$ when $u,y$ are fresh.
For the padding converse, eliminate the paired formula, supply its
zero-gate theorem, and Use its clause at $1$. The gate is essential;
an unused $D$ without that gate would make a different formula.

Define the selector by recursion on the maximum binder depth:
\[
\begin{aligned}
M_v(s=t,r=w)&:=(C(v,s,r)=C(v,t,w)),\\
M_v(D_{x,y}A,D_{x,y}B)&:=D_{x,y}M_v(A,B),\\
M_v(E,D_{x,y}B)&:=D_{x,y}M_v(G_yE,B),\\
M_v(D_{x,y}A,E)&:=D_{x,y}M_v(A,G_yE),
\end{aligned}
\]
where $E$ is an equation and paired binders are aligned and fresh.
Padding and matrix congruence prove
$M_0(A,B)\dashv\vdash A$ and $M_{Sz}(A,B)\dashv\vdash B$.
Let $1=S0$, $2=S1$, and
$p(z)=\si(z,0;u.C(u,1,0))$; only the computations
$p(1)=1$ and $p(2)=0$ are needed here.

With all new binders fresh and the first coordinate of $q$ unused,
the full logical macro inventory is
\[
\begin{aligned}
\bot_*&:=(0=1),\\
A\Rightarrow_*B&:=q_yM_y(A,B),\\
A\wedge_*B&:=q_yG_yM_{p(y)}(A,B),\\
\exists_*x\,A(x)&:=D_{x,y}G_yA(x),\\
A\vee_*B&:=\exists_*v\,M_v(A,B),\\
\forall_*x\,A(x)&:=A(0)\wedge_*q_yA(y).
\end{aligned}
\]
These are finite formula translations, not additional kernel rules.
Implication introduction uses case introduction; its elimination Uses the
opened clause at $1$. Existential introduction uses the selected witness
and the successor gate; its elimination supplies the true zero gate and
Uses at $1$. Conjunction projections use the true zero gate and the two
instances $1,2$. Its introduction follows by numerical cases on the
selector control. Disjunction uses existential witnesses $0,1$; numerical
cases on its control give the elimination rule. The two universal rules
follow from the $q$ interfaces and the conjunction rules.

Finally, under $0=1$, transport through $C(v,s,t)$ gives $s=t$ for
arbitrary terms. External induction on formulas with case introduction
then gives every formula, proving falsity elimination. The inherited
development in Appendix~\ref{sec:choice-introduction} gives these
derivations in detail; the normalized rules above supply each ingredient.
No excluded middle for quantified formulas is assumed.

\subsection{The internal arithmetic bridge}

The inherited arithmetic development is included in the appendices, rather
than replacing an internal interpretation argument by computations on
standard numerals. Its essential dependency order is:
\begin{enumerate}
\item Derive a predecessor $P$ with $P0=0$ and $P(Sn)=n$ from scalar
iteration, equation induction, and numerical cases.
\item Build truncated subtraction and a faithful term characteristic of
actual equality; its reflection theorem holds at arbitrary model elements.
\item Compile every fixed quantifier-free intuitionistic formula to a
characteristic equation, obtaining the full quantifier-free induction scheme.
\item Construct pairing and projections by bounded search, then compile
every fixed primitive recursion to scalar iteration.
\end{enumerate}
The predecessor is established before assuming a pairing function or a
faithful equality detector. This ordering is necessary to avoid circularity.

For pairing $\pi$ with projections $L,R$, the recursion
\[
f(0,\vec p)=g(\vec p),\qquad
f(Sn,\vec p)=h(n,f(n,\vec p),\vec p)
\]
is represented by
\[
\begin{split}
 V(n,\vec p)&=I\bigl(n,\pi(0,g^*(\vec p));
 u.\pi(SL(u),h^*(L(u),R(u),\vec p))\bigr),\\
 f^*(n,\vec p)&=R(V(n,\vec p)).
\end{split}
\]
Equation induction proves $L(V(n,\vec p))=n$ and the defining
equations for $f^*$. Each star is a finite syntactic compilation of a
fixed description. No total numerical evaluator for all primitive-recursive
descriptions is introduced.

\subsection{Why the result is a bi-interpretation}

Interpret scalar iteration in ordinary iPRA by primitive recursion whose
step ignores the recursion counter. Interpret each iPRA function back into
$T_I$ by the scalar compiler. The domain formula in each direction is
$x=x$, and equality is unchanged. Uniqueness of each fixed recursion follows
from equation induction. External induction on the finite description then
proves, internally, that each doubly translated function equals its original.
This also supplies coherence when substitution changes an update description.

The numerical identity graph $x=y$ is therefore internally a total
bijection preserving every operation in both composites. Formula round
trips follow from these numerical equalities, the derived intuitionistic
logical rules, and the native/$D$ equivalences of
Theorem~\ref{thm:roundtrips}. In particular the paired formula has the
expected internal equivalence with
\[
\exists_*x\bigl(B(x,0)\Rightarrow_*\forall_*yB(x,y)\bigr).
\]
The forward direction opens its clause and proves all instances under the
base; the reverse direction uses the given universal in case introduction.
This proves Theorem~\ref{thm:main}. The translation also identifies
$0=S0$ with ordinary falsity, yielding equiconsistency.

The pinned Lean endpoints \file{native_iff_pra},
\file{pra_iff_native}, \file{native_roundTrip},
\file{pra_native_roundTrip}, and \file{native_inconsistent_iff}
document the inherited bridge~\cite{bridge}. The new clause/single-Eval
comparison is the written proof above. These existing source files were
inspected and are included as references; they were not rebuilt in this run.

\section{Complete implementation and verification}

\subsection{Public contract and trusted boundary}

The selected checker is a single Python standard-library module. Its input
is one JSON object with exactly \file{goal} and \file{proof}. The goal
must be a closed normal Eq/$D$ formula. Its successful public result is
exit status zero and exactly \file{ACCEPT} followed by a newline. Invalid
certificates print \file{REJECT} and exit with status one.

The numeral $0$ denotes zero; positive integers are local bound indices,
starting at 1; negative integers are fresh parameter levels. At ambient
scope $n$, the available free parameters are $-1,\ldots,-n$. Successor
is \file{[0,t]}, iteration is \file{[n,a,h]}, equality is
\file{[s,t]}, and $D$ is \file{[B]}. An iterator binds one index only
in its third entry; $D$ binds two, with $y$ at index 1 and $x$ at index 2.
Term and formula validation use the expected sort, so an equality pair
cannot silently stand for a successor term.

The five proof patterns are represented as follows. A validated reflexive
goal closes immediately; its proof payload is unused, and may be JSON null.
Every control uses the same convention. Other nodes use these packets:
\begin{center}
\begin{tabular}{@{}ll@{}}
\toprule
Rule & Packet\\
\midrule
Eval & \file{['v',n,a,h,t,major,minor]}\\
Use & \file{['u',index,t,premise]}\\
Elim & \file{['e',body,major,branch]}\\
Intro & \file{['i',witness,branch]}\\
\bottomrule
\end{tabular}
\end{center}
These are readable Python-style displays of JSON arrays; actual files use
double-quoted strings. Eval constructs its predecessor and successor
equations from the supplied counter, seed, and update. In a goal $r=w$,
it checks $w=h\langle t/u\rangle$, then verifies both prescribed premises.
The constructed successor premise validates all the counter, seed, and
update data, including data that a zero predecessor would discard. The
major premise validates $t$. This avoids redundant independent traversals
without weakening the validation invariant.

The opening traversal does not shift its inserted value because every
replacement is a depth-zero term at the current ambient scope: only $0$
and negative free levels occur outside its own binders. It still tracks
the binders traversed inside the target. Elim validates the parent goal
before admitting a fresh level, preventing eigenvariable escape; Intro
validates its witness before opening its fresh step variable.

All syntax validation, normalizing opening, branch scope, clause handling,
JSON parsing, and public output are included in the size measurements.
The proof compiler, tests, inherited sources, and operating-system/Python
runtime are not part of the standalone kernel. They are reported separately.
Acceptance does not depend on Python \file{assert}, so optimization mode
does not remove proof checks.

\subsection{Executable evidence}

The packaged reproduction run passed all of the following:
\begin{itemize}
\item 66 raw-to-normal certificate translations, each checked in the
raw CE5 verifier, the forward single-Eval verifier, the selected
checker, the matched controls, and an independent archived normalized
Eq/$Q$/$D$ verifier after the reverse translation.
\item 70 accepted certificates and 20 adversarial rejected certificates.
\item 180 public-interface executions, once in ordinary Python and once
with \file{python -O} for each accepted or rejected certificate.
\end{itemize}
The 66 source certificates comprise 32 original CE5 examples, 32 distinct
round-tripped proof certificates for those examples, and two explicit
witness-normalization repairs. They are not 66 independent mathematical
phenomena. Four additional positive cases exercise a genuinely used
induction hypothesis and three closure targets. Rejections include false
equations, malformed or incorrectly scoped terms, Boolean or floating-point
indices, wrong evaluation premises, illegal witnesses, and an invalid
change from progressive to case introduction.

Two independent reviews examined the mathematical comparison and the
implementation. The implementation review confirmed the validation-through-premises invariant,
the scope discipline, and the harmless unused payload at reflexive leaves.
The packaged \file{evidence/verification.json} records the reproducible
run; \file{evidence/implementation_audit.txt} records the review boundary.

The forward compiler preserves the end formula up to zero normalization.
It explicitly repairs a case-introduction witness that normalizes to $0$,
using the surviving alias $Z$. A reverse compiler checks produced proofs
against an independently supplied normalized native checker with Eq/$Q$/$D$
formulas. These transformations are untrusted conveniences: the new checker
validates their output independently.

The finite checks establish implementation evidence for the supplied
examples and adversarial cases. They do not prove the universal mathematical
metatheorems or certify the Python implementation in Lean. Proof-tree
translation, checker correspondence, and the new complete bi-interpretation
remain written mathematical arguments, with the stated independent reviews.

\subsection{Reproduction and limits}

Run \file{python3 reproduce.py} from the extracted package. It reruns the
portable certificate checks and recreates the source metrics and normalized
gzip files. Run \file{pdflatex -interaction=nonstopmode -halt-on-error report.tex}
twice to rebuild the PDF. The delivered \file{report.tex} is standalone;
it embeds the checker and the attributed arithmetic appendix.

The checker is a terminating recursive algorithm on finite certificates
and terms, abstracting from finite machine resources. Python recursion or
memory limits can reject very large valid certificates, as for the matched
controls. Only zero iteration is normalized, so no full primitive-recursive
program evaluator is trusted. The new implementation has not been proved
extensionally correct by a proof assistant, and no global optimality claim
is made.

\appendix
\clearpage
\section*{Attribution and inherited foundations}
\addcontentsline{toc}{section}{Attribution and inherited foundations}

The following five mathematical components are reproduced from the
repository's delivered reference development, preserved in the archive
under \file{reference/} with its MIT No Attribution license and provenance.
They establish the inherited choice basis, ordinary bridge, scalar
predecessor, equality reflection, and primitive-recursion compiler. Their
raw syntax and larger rule inventory belong to that reference development;
Sections~\ref{sec:specification}--\ref{sec:interpretation} above prove the
comparison with the present normal five-pattern kernel. They are supplied
to make the arithmetic part of the interpretation inspectable without
network access, not as new contributions of this report.

% New mathematical result: complementary introduction choices.
% Requires amsmath/amssymb/amsthm and theorem, lemma, proposition environments.
% The inherited arithmetic-to-iPRA bridge must be supplied separately.

\section{Complementary introductions with one local assumption}
\label{sec:choice-introduction}

The old paired introduction retains both the zero instance and an
induction instance when its matrix is an equation. Both together are
unnecessary. A certificate may choose either one, provided the induction
choice is restricted to a literal equation. This changes neither the
formula language nor the meaning of a formula. The omitted combination
is recovered by a short nested introduction--elimination construction.

\subsection{Complete native syntax and rule inventory}

The scalar numerical language is
\[
 t::=x\mid0\mid St\mid I(n,a;u.h).
\]
Only $h$ binds the state variable $u$; $n$ and $a$ do not.
The update may contain fixed numerical parameters. The computation
axioms in every retained context are
\[
 I(0,a;u.h)=a,\qquad
 I(Sn,a;u.h)=h[I(n,a;u.h)/u].
 \tag{I0, IS}\label{eq:choice-computations}
\]
The argument also applies unchanged to the full fixed-description PR
language with its projection, composition, and recursion equations.
No numerical variable ranges over function descriptions, and no uniform
evaluator is postulated.

Native formulas and reusable context clauses are
\[
 A::=(s=t)\mid\mathsf D_{x,y}A,\qquad
 \langle A;y.B(y)\rangle,\quad y\notin\operatorname{FV}(A).
\]
Bound variables are renamed before capture-free substitution.
The two variables $x,y$ in a paired formula are distinct.
An ordinary assumption $F$ is the constant clause
$\langle H;v.F\rangle$, where $v$ is absent from $F$ and $H$ is the
fixed computation theorem $I(0,0;u.0)=0$. Thus ordinary hypothesis
use is a clause Use preceded by this computation proof.

There are computation axioms and exactly the following inference
patterns:
\[
\frac{\Gamma\vdash A}{\Gamma\vdash B(t)}
\quad\mathsf{Use}
\quad(\langle A;y.B(y)\rangle\in\Gamma),
\]
\[
\frac{\Gamma\vdash\mathsf D_{x,y}B(x,y)\qquad
      \Gamma,\langle B(a,0);y.B(a,y)\rangle\vdash F}
     {\Gamma\vdash F}
\quad\mathsf{DE},
\qquad
\frac{\Gamma\vdash s=t\qquad\Gamma\vdash F(s)}
     {\Gamma\vdash F(t)}
\quad\mathsf{Tr},
\]
\[
\frac{\Gamma,B(t,\eta_\varepsilon(z))\vdash B(t,Sz)}
     {\Gamma\vdash\mathsf D_{x,y}B(x,y)}
\quad\mathsf{DI}_{\varepsilon},
\qquad
\eta_0(z)=0,\quad\eta_1(z)=z.
\tag{Choice DI}\label{eq:choice-di}
\]
Here $\varepsilon$ is a certificate choice in $\{0,1\}$, and
$\varepsilon=1$ is permitted only when $B$ has a literal equation
root. An equation beneath a binder is not a literal equation.
Call $\varepsilon=0$ the \emph{case choice} and $\varepsilon=1$
the \emph{progress choice}. Both retain exactly one ordinary
formula assumption. Compound matrices have only the case choice.

In DI the eigenparameter $z$ is fresh for $\Gamma$, the witness
$t$, and every fixed parameter of $B$. In DE the eigenparameter
$a$ is fresh for $\Gamma$, the conclusion $F$, and every fixed
parameter of $B$. In transport, $F(v)$ is any fixed native
formula template. Displayed local assumptions and clauses are
discharged. There are no primitive reflexivity, coverage,
separation, ordinary logical, or additional induction rules.

We call this calculus $\mathcal C$. The original comparison
calculus $\mathcal N$ has identical syntax, computation,
Use, DE, and transport, but its introduction is
\[
\frac{\Gamma,B(t,0),[B(t,z)]_{\mathrm{at}}\vdash B(t,Sz)}
     {\Gamma\vdash\mathsf D_{x,y}B(x,y)}.
\tag{Old DI}\label{eq:choice-old-di}
\]
The bracket supplies $B(t,z)$ exactly at a literal equation.

\subsection{Exact equivalence and the unused-witness construction}

\begin{lemma}[Structural admissibility]
\label{lem:choice-structural}
Weakening, exchange, contraction, renaming, ordinary-assumption cut,
and capture-free numerical substitution are admissible in $\mathcal C$.
\end{lemma}
\begin{proof}
Induct on derivations. Rename local eigenparameters before inserting
proofs or terms. Numerical substitution preserves a formula's root
constructor, so an allowed progress choice remains allowed.
The selected choice does not change under substitution or cut.
For ordinary-assumption cut, replace each use of its constant clause
by the supplied proof. Computation schemas are closed under
capture-free numerical substitution. These facts verify every case.
\end{proof}

\begin{theorem}[Identity-language equivalence]
\label{thm:choice-equivalence}
For every retained clause context $\Gamma$ and native formula $F$,
\[
 \Gamma\vdash_{\mathcal C}F
 \quad\Longleftrightarrow\quad
 \Gamma\vdash_{\mathcal N}F.
\]
The transformations are effective on finite derivations and leave
all numerical and formula syntax unchanged.
\end{theorem}
\begin{proof}
Each new introduction is an instance of the old one after weakening
by the omitted local assumption. For the case choice at an equation,
add the induction instance. For the progress choice, add the zero
instance. Compound case introduction is already the old rule.
Thus every new derivation translates to an old one.

For the reverse direction, transform an old derivation recursively.
Only its atomic DI case needs a construction. Suppose its matrix is
the equation $E(x,y)$ and its transformed premise is
\[
 \pi:\quad\Gamma,E(t,0),E(t,z)\vdash E(t,Sz).
\]
The outer new inference will be case DI with witness $t$. It is
therefore enough to prove $E(t,Sz)$ under
$\Gamma,E(t,0)$.

Rename the eigenvariable in $\pi$ to a fresh $v$. Choose a new
variable $u$ absent from $E(t,y)$ and all retained data. Under
$\Gamma,E(t,0)$, progress DI with witness zero introduces
\[
             \mathsf D_{u,y}E(t,y)
\]
from exactly the renamed premise $\pi[v/z]$.
The bound coordinate $u$ is deliberately unused.

Apply DE to this paired formula. Because $u$ is unused, the opened
clause is exactly
\[
               \langle E(t,0);y.E(t,y)\rangle;
\]
it contains no new existential witness. The retained ordinary
assumption $E(t,0)$ supplies its antecedent. Use at $Sz$ gives
$E(t,Sz)$, and DE closes. Finally the outer case DI introduces
the original $\mathsf D_{x,y}E(x,y)$ with witness $t$.
The original premise proof $\pi$ has been inserted once.

An old compound introduction already is the new case choice.
All other rule cases are unchanged. This completes the induction
on the original derivation.
\end{proof}

The unused inner coordinate is essential. Introducing
$\mathsf D_{x,y}E(x,y)$ internally and then eliminating it would
give an arbitrary fresh witness, not necessarily the original
$t$. The construction instead keeps $t$ as a fixed parameter
of $E(t,y)$ and binds only the unused coordinate $u$.

\begin{proposition}[Linear bound on inference nodes]
\label{prop:choice-node-bound}
Count each computation axiom, Use, DI, DE, and transport as one
inference node, expanding ordinary-assumption use as a computation
of $H$ followed by Use. If an old derivation has $N$ nodes and
$r$ atomic DI nodes, the displayed transformation has exactly
\[
                   N+5r\leq 6N
\]
nodes.
\end{proposition}
\begin{proof}
Each atomic old DI node is replaced by six nodes: outer case DI,
inner progress DI, DE, Use of the opened clause, Use of the
retained base assumption, and the computation of $H$ needed for
that last ordinary-assumption use. Its old premise subtree
occurs once. All other inference nodes are copied once.
\end{proof}

This counts inference-tree nodes. It is not a bound on serialized
JSON bytes, formula sizes, checker runtime, or the number of
consumed packet cells when hypothesis closure is implicit.
In the formula-context Q/D presentation, the corresponding
replacement has outer case DI, inner progress DI, DE, QE,
and two Hyp nodes, giving the same six-node count.

\subsection{Induction and numerical cases separate cleanly}

\begin{lemma}[Contextual equation induction]
\label{lem:choice-ei}
The progress choice, DE, and Use derive
\[
\frac{\Gamma\vdash E(0)\qquad
      \Gamma,E(z)\vdash E(Sz)}
     {\Gamma\vdash E(t)}
\tag{EI}\label{eq:choice-ei}
\]
for every equation $E$, with the usual freshness condition.
\end{lemma}
\begin{proof}
Choose an unused witness coordinate $u$. Progress DI introduces
$\mathsf D_{u,y}E(y)$ from the given step. DE opens the clause
$\langle E(0);y.E(y)\rangle$. Supply its antecedent using the given
base proof and use the clause at $t$. Close DE.
\end{proof}

\begin{lemma}[Cases, equality, and binder congruence]
\label{lem:choice-preinduction}
The case choice derives, for arbitrary native formulas,
\[
\frac{\Gamma\vdash F(0)\qquad\Gamma\vdash F(Sz)}
     {\Gamma\vdash F(t)}
\tag{NC}\label{eq:choice-nc}
\]
with the usual freshness condition. Equality reflexivity,
symmetry, transitivity, numerical congruence, and uniform
matrix interderivability under $\mathsf D$ are derivable
without using the progress choice.
\end{lemma}
\begin{proof}
For cases, use case DI with an unused witness coordinate and
matrix $F(y)$, ignoring its retained base in the successor
premise. DE opens $\langle F(0);y.F(y)\rangle$; the given zero
proof and Use at $t$ finish.

For reflexivity put $r=I(0,t;u.0)$. Use the computation equation
$r=t$ as both transport premises with motive $v=t$ to obtain
$t=t$. The usual explicit equality motives give symmetry,
transitivity, and numerical congruence.

For binder congruence, open the old paired formula at a fresh
witness $a$ and introduce the new one by case DI at $a$.
Convert its retained base to the old base, use the opened
old clause at the successor, and convert back. Ordinary cut
closes the construction, and the reverse conversion gives
the converse. No induction is used.
\end{proof}

The primitive progress choice is still restricted to equations.
Allowing it at arbitrary compound matrices would, by the
argument for \eqref{eq:choice-ei}, supply induction for
arbitrary formulas and change the intended theory.

\subsection{The entire ordinary logical interface}
\label{subsec:choice-logic}

Here is a predecessor-free reconstruction of the ordinary logical
operations. It makes explicit that the case choice supplies the
logical interface independently of progress induction.

Let $C(v,a,b):=I(v,a;u.b)$ with $u$ unused. Computation gives
$C(0,a,b)=a$ and $C(Sv,a,b)=b$. Define the gate on fixed finite
formula syntax by
\[
 G_v(s=t):=(C(v,0,s)=C(v,0,t)),\qquad
 G_v(\mathsf D_{x,y}B):=\mathsf D_{x,y}G_vB,
\]
renaming bound variables away from the control. Equality,
binder congruence, and external structural induction give
\[
 \vdash G_0A,\qquad G_{Sv}A\dashv\vdash A,\qquad
 A\dashv\vdash\mathsf D_{u,y}G_yA
 \quad(u,y\notin\operatorname{FV}(A)).
\tag{Gate and padding}\label{eq:choice-gate}
\]
For padding, case DI proves the forward direction. DE, the
zero-gate theorem, and Use at $1$ prove the reverse direction.

Define a selector between two arbitrary formulas by
\[
 M_v(s=t,u=w):=(C(v,s,u)=C(v,t,w))
\]
and, with aligned fresh paired variables, by
\[
\begin{aligned}
 M_v(\mathsf D_{x,y}A,\mathsf D_{x,y}B)
   &:=\mathsf D_{x,y}M_v(A,B),\\
 M_v(E,\mathsf D_{x,y}B)
   &:=\mathsf D_{x,y}M_v(G_yE,B),\\
 M_v(\mathsf D_{x,y}A,E)
   &:=\mathsf D_{x,y}M_v(A,G_yE),
\end{aligned}
\]
where $E$ is an equation. Every call decreases maximum input
binder depth. Padding and congruence prove
$M_0(A,B)\dashv\vdash A$ and
$M_{Sv}(A,B)\dashv\vdash B$.

Put $\delta(v)=C(v,1,0)$ and $p(z)=I(z,0;u.\delta(u))$.
Only the fixed computations $p(1)=1$ and $p(2)=0$ are needed.
The full logical macro inventory is
\[
\begin{aligned}
 A\Rightarrow_*B&:=\mathsf D_{u,z}M_z(A,B),\\
 A\wedge_*B&:=\mathsf D_{u,z}G_zM_{p(z)}(A,B),\\
 \exists_*x\,A(x)&:=\mathsf D_{x,y}G_yA(x),\\
 A\vee_*B&:=\exists_*x\,M_x(A,B),\\
 \forall_*x\,A(x)&:=
        A(0)\wedge_*\mathsf D_{u,z}G_zA(z),\\
 \bot_*&:=(0=1).
\end{aligned}
\tag{Logical macros}\label{eq:choice-macros}
\]
Newly bound variables are fresh; $u$ is unused.

Implication introduction obtains its assumption from the zero
selector and supplies the consequent at successors, then uses
case DI. Elimination opens its clause and uses it at $1$.
Existential introduction uses its witness and the successor
gate; elimination opens at a fresh witness, supplies the
zero-gate theorem, and uses the clause at $1$.

Disjunction introduction uses witnesses $0$ and $1$.
For elimination, discharge the two branch proofs as
$A\Rightarrow_*F$ and $B\Rightarrow_*F$. NC applied to
$M_z(A,B)\Rightarrow_*F$ gives this implication at every
control: the two endpoint cases follow from the branch
implications. Existential elimination provides $M_a(A,B)$
and the implication gives $F$.

Given $A$ and $B$, selector endpoints and NC give $M_v(A,B)$
at an arbitrary control. Substitute $v=p(Sz)$ and use the
successor gate for conjunction introduction. For projections,
open its zero-gated clause and specialize at $1$ and $2$;
the two fixed computations give $B$ and $A$, respectively.

These rules prove existential numerical coverage:
\[
     (x=0)\vee_*\exists_*v\,(x=Sv).
\tag{Coverage}\label{eq:choice-coverage}
\]
Use NC. Zero takes the left injection of reflexivity, and a
successor takes the right injection with the evident witness.
Coverage gives equality cases even when the retained context
contains $x$; its elimination adds $x=0$ or $x=Sv$ rather
than silently changing that retained context.

Universal introduction follows from a fresh-parameter proof
by substituting zero and successors and forming the two
conjuncts. For elimination, take its $A(0)$ conjunct and
open the paired conjunct. Its base is a zero-gate theorem.
Coverage of the target chooses the zero instance from
$A(0)$ or a successor instance from the opened clause
and gate removal.

Finally, under $0=1$, computation and transport give
$s=C(0,s,t)=C(1,s,t)=t$ for arbitrary numerical terms.
Every equation follows, and external induction on formula
syntax with case DI gives every native formula. This
supplies falsity elimination. No excluded middle for
quantified formulas has entered the construction.

\subsection{Whole-domain realizations and both formula composites}

For a standard first-order reading, let $T_{\rm eq}$ have
intuitionistic equality logic, the same numerical signature
and computation equations, numerical coverage, separation
$(0=1)\to\bot$, and contextual equation induction.
Its interpretation of the native syntax is
\[
 (s=t)^\circ=(s=t),\qquad
 (\mathsf D_{x,y}B)^\circ
   =\exists x\bigl(B^\circ(x,0)\to\forall y B^\circ(x,y)\bigr),
\]
\[
 \langle A;y.B\rangle^\circ
       =A^\circ\to\forall y B^\circ.
\]
The domain is the entire numerical domain and interpreted
equality is actual equality.

Case DI is sound in this realization by numerical coverage.
Progress DI is sound by contextual equation induction under
the conditional base. Use, DE, transport, and computation
have their displayed ordinary meanings. Conversely, the
native macros derive every ordinary logical rule,
\eqref{eq:choice-ei} derives equation induction, and
\eqref{eq:choice-coverage} derives coverage. Separation
translates to the identity implication
$(0=1)\Rightarrow_*(0=1)$.

For the essential paired formula composite,
\[
 \mathsf D_{x,y}B
 \dashv\vdash
 \exists_*x\bigl(B(x,0)\Rightarrow_*\forall_*y\,B(x,y)\bigr).
\tag{Paired composite}\label{eq:choice-paired-composite}
\]
Forwards, DE opens its clause at a fresh $a$; under $B(a,0)$
that clause supplies $B(a,y)$ at a fresh arbitrary $y$.
Introduce the universal, implication, and existential
macros and close DE. Backwards, eliminate the existential
at a fresh $a$. Introduce the paired formula with witness
$a$ by case DI: the retained base, supplied implication,
and universal elimination at $Sz$ give its successor premise.

The ordinary interpretations of the six logical macros
have their ordinary round trips. For conjunction, its true
base and the instances $1,2$ recover both operands;
the converse uses selector cases. The universal has all
successor instances in its paired conjunct and zero in
its first conjunct; coverage gives exactly the ordinary
universal. Implication, disjunction, and existential
quantification use their selector or gate endpoints.
Falsity uses the explicit ordinary separation bridge.
External induction on finite formulas, together with
\eqref{eq:choice-paired-composite}, gives both formula
composites. Clauses correspond in both directions to
$A\Rightarrow_*\forall_*yB$ by the derived logical rules.

The numerical composites are literal identities with
internally provable graph $x=y$. No standard-numeral
restriction, retract, quotient, or equality reflection
assumption is being inserted at this stage. The separately
supplied arithmetic bridge identifies $T_{\rm eq}$ with
full iPRA; the present change preserves that bridge
unchanged because Theorem~\ref{thm:choice-equivalence}
preserves every native sequent.

\subsection{What the theorem does and does not claim}

The new inference interface retains one ordinary assumption
at every paired introduction. It preserves the entire
consequence relation, not just equations or closed
theorems. Consequently the established full-iPRA
interpretations and equiconsistency transfer unchanged.

This result is different from retaining only progress
introduction and strengthening closure elimination.
Here the original elimination rule is retained, and
the case and progress introductions coexist. It is also
different from allowing arbitrary compound progress,
which would supply stronger induction.

A complete-checker comparison must still count the
certificate choice, its equation-root restriction,
substitution, context handling, computation, and the
public input/output interface. The mathematical theorem
alone is not a byte-count or minimality theorem.
The inference-node bound is distinct from source-size
measurements and from serialized certificate size.
The proof is a written universal proof; finite regression
tests do not constitute a Lean/Rocq formalization of it.


\section{Arithmetic strength and both ordinary interpretation composites}
\label{sec:ordinary-bridge}

We spell out the arithmetic bridge to avoid identifying standard numerical
agreement with bi-interpretability. First elaborate every fixed iterator
into an ordinary function symbol. At stage zero the symbols are $0,S$.
Given a fixed earlier-stage term $h(u,\vec p)$, introduce
$i_h(n,a,\vec p)$ with the two equations
\[
i_h(0,a,\vec p)=a,\qquad
i_h(Sn,a,\vec p)=h(i_h(n,a,\vec p),\vec p).
\]
Only finite descriptions occur; the domain contains no code for a
universal evaluator. Formal parameters are distinct when a description
is declared, while actual arguments can coincide.

Let $T_I$ be ordinary intuitionistic first-order logic with actual
equality and genuine falsity, these staged equations, contextual
induction for literal equations, and
\[
\forall x\,(x=0\vee\exists y\,x=Sy),
\qquad (0=S0)\to\bot.
\]
The second sentence is the falsity bridge. It must be stated: positive
numerical equations by themselves admit a singleton structure. Native
explosion from $0=S0$ and a translation of ordinary falsity to that
equation do not authorize silently dropping the bridge from $T_I$.

The native-to-ordinary translation is
\[
(s=t)^\circ=(s=t),\qquad
(\mathsf D_{x,y}B)^\circ=
\exists x\bigl(B^\circ(x,0)\to\forall y B^\circ(x,y)\bigr).
\]
A clause $\langle A;y.B\rangle$ becomes $A^\circ\to\forall yB^\circ$.
For the auxiliary closure used by the compact checkers,
$(\mathsf Q_yB)^\circ=B^\circ(0)\to\forall yB^\circ$.
The interpretation domain is all elements, defined by $x=x$;
equality is actual equality.

Case DI is sound because numerical coverage combines its retained base
with its successor proof. Progress DI is sound by literal-equation
induction under a temporarily assumed base. Use, DE and transport are
ordinary logical inferences under the translations. Conversely the
derived logical operations, equation induction, cases and explosion in
Section~\ref{sec:choice-introduction} recover the rules and axioms of
$T_I$. An auxiliary closure is interderivable with
$\mathsf D_{v,y}B(y)$ for unused $v$: in one direction use case DI with
the closure and its base; in the other eliminate the paired formula
and return its supplied closure. Thus formulas and context clauses have
the stated meanings in both presentations.

\subsection{The arithmetic dependency order}

The scalar update has no recursion-counter argument. We do not conceal
one in an encoding. The inherited internal proofs, reproduced in the
appendices, establish the required arithmetic in this order:
\begin{enumerate}
\item Using equation induction and numerical cases, construct a fixed
scalar predecessor term and prove $P0=0$ and $P(Sn)=n$.
Existential coverage then gives $x=0\vee x=S(Px)$: in the successor
case $x=Sy$, apply $P$ to obtain $Px=y$ and substitute.
\item Construct truncated subtraction and a numerical characteristic
of actual equality. Prove its reflection on the entire original domain.
\item Compile each fixed quantifier-free formula into its characteristic
equation and recover every contextual quantifier-free induction instance.
\item With that induction available, construct pairing and projections
by a bounded square-root search and compile every fixed general
primitive-recursive description into a scalar iterator term.
\end{enumerate}
No pairing or faithful equality detector is assumed to obtain
predecessor. No appeal to standard-model correctness replaces an
internal equality. Every induction fixes its passive parameters;
diagonal substitution occurs only into a previously proved identity.

For a target recursion
\[
f(0,\vec p)=g(\vec p),\qquad
f(Sn,\vec p)=h(n,f(n,\vec p),\vec p),
\]
let $\pi,L,H$ be the proved pairing and projections. The compiler uses
\[
V(n,\vec p)=I\bigl(n,\pi(0,g^*(\vec p));
u.\pi(SL(u),h^*(L(u),H(u),\vec p))\bigr),
\qquad f^*(n,\vec p)=H(V(n,\vec p)).
\]
Equation induction proves $L(V(n,\vec p))=n$; the two projection
identities then prove the target defining equations. The construction
is external recursion on each finite PR description, not an internal
quantification over arbitrary functions.

\begin{theorem}[Whole-domain bi-interpretability and equiconsistency]
The ordinary realization $T_I$ and full first-order iPRA are
bi-interpretable. In both theories $x=y$ defines the isomorphism from
the original domain to its double interpretation. The native Choice-DI
calculus is therefore a faithful proof presentation of a theory
bi-interpretable and equiconsistent with iPRA.
\end{theorem}
\begin{proof}
Interpret an iterator in iPRA by primitive recursion that ignores its
counter argument in the update. iPRA proves its equations, literal
equation induction, coverage and separation. In the reverse direction
use the scalar compiler. The predecessor and equality-reflection proofs
supply the successor axioms and every quantifier-free induction instance.
All ordinary logical rules are preserved.

For two iterators with the same seed and provably equal updates,
equation induction on the equality of their values proves uniqueness at
every counter, with all parameters fixed. The corresponding argument
applies to ordinary primitive recursion. Induct externally on a finite
description: its double translation satisfies its original defining
equations, so internal uniqueness identifies the two functions. This
also proves coherence when substitution changes the staged syntactic
description of an update.

The domain and equality are unchanged; hence the identity graph is
internally total, bijective and preserves every interpreted operation.
The displayed ordinary meanings of the logical macros and the paired
binder give both formula composites by external induction on finite
formula syntax. Clauses are similarly equivalent to their conditional
universal meanings. The explicit falsity bridge preserves contradiction
in both directions. Effective translations of finite proofs give
equiconsistency, without a claim that either theory proves its own
consistency.
\end{proof}


% Adapted with attribution from the pinned unparenthesized one-hole archive.
\section{Internal scalar predecessor bootstrap}
\label{sec:hybrid-bootstrap}

This scalar bootstrap is inherited from the supplied scalar development~\cite{erasure} and is reproduced to keep the scalar checker self-contained. It is not a new contribution of this paper. The proof uses a factor map fixing odd states and a saturated-iteration construction of predecessor.

The available arithmetic terms are $0$, $S$, and
\[
 I(n,a;u.h),\qquad
 I(0,a;u.h)=a,\qquad
 I(Sn,a;u.h)=h[I(n,a;u.h)/u].
\]
The update binds only $u$.  All other variables in it are fixed numerical
parameters.  We use equality reasoning, the already reconstructed
zero/successor cases, and induction on a single literal equation.
Every induction below keeps its other parameters fixed.  Predecessor,
successor injectivity, pairing, and induction on arbitrary
quantifier-free formulas are not assumptions of the argument.

\subsection{Elementary terms and two finite cycles}

Write $1=S0$ and use larger numerals as successor abbreviations.  Define
\[
 C(n,a,b)=I(n,a;u.b),\qquad
 Z(n)=C(n,1,0),\qquad
 a+b=I(b,a;u.Su).
\]
The update variable of $C$ is absent from its body.  Consequently
\[
 C(0,a,b)=a,\qquad C(Sn,a,b)=b.
\]
Numerical cases show that $Z(n)$ is $0$ or $1$.
The usual commutative-monoid equations for addition follow by literal
equation induction.  For example, $0+b=b$ and $Sa+b=S(a+b)$ are
proved by induction on $b$; these yield commutativity by a further
equation induction, and associativity follows by induction on the last
summand.  We use these equations to reorder finite sums.  The notation
$2a$ means $a+a$ and $3a$ means $a+a+a$; no multiplication operation
is needed for this bootstrap.

Put
\[
 N(n)=I(n,0;u.Z(u)),\qquad
 j(u)=Z(u)+2N(u),\qquad
 M(n)=I(n,0;u.j(u)).
\]
The range of $N$ is contained in $\{0,1\}$: its zero value is $0$,
and its successor value is a value of $Z$.  This uses numerical cases
on its counter, not induction on a disjunction.  Hence
\begin{equation}
 N(n)+Z(N(n))=1,\qquad
 Z(Z(N(n)))=N(n),\qquad N(n+2)=N(n).
 \label{eq:hybrid-parity}
\end{equation}
Here the last equation follows from two successor computations and
the preceding equation.  Literal-equation inductions give
\begin{equation}
 N(2n)=0,\qquad N(2n+1)=1,\qquad N(3n)=N(n).
 \label{eq:hybrid-parity-scaling}
\end{equation}
For the first equation the successor step adds two to the argument of
$N$; the second follows by one computation.  In the third, increasing
$n$ adds three to the left argument and therefore applies $Z$ once to
both sides of the induction equation.

Finite computations now give
\[
 j(0)=1,\qquad j(1)=2,\qquad j(2)=0.
\]
The proof below needs just these three values and
$M(Sn)=j(M(n))$; it does not assume a division algorithm or a theorem
characterizing $M$ as remainder modulo three.

\subsection{A ramp with a global scaling identity}

Define
\[
 G(n)=I(n,0;u.u+M(u)+1).
\]
The central invariant is
\begin{equation}
 M(G(n))=N(n).
 \label{eq:hybrid-ramp-invariant}
\end{equation}
Its zero case is immediate.  For its successor step, retain the
single induction equation and split the Boolean value $N(n)$.
If it is $0$, the next $G$ state is $G(n)+1$; its $M$ value is
$j(0)=1$, equal to $N(Sn)$.  If it is $1$, the next $G$ state
is $G(n)+2$; two computations of $M$ change its value from $1$
to $2$ and then to $0$, again equal to $N(Sn)$.  This proves
the invariant by literal-equation induction.

It follows that
\[
 G(Sn)=G(n)+N(n)+1,\qquad G(n+2)=G(n)+3.
\]
Equation induction therefore yields
\begin{equation}
 G(2n)=3n.
 \label{eq:hybrid-ramp-double}
\end{equation}
A separate literal-equation induction proves
\begin{equation}
 G(3n)=3G(n)+N(n).
 \label{eq:hybrid-ramp-triple}
\end{equation}
For completeness, its successor calculation is
\begin{align*}
 G(3Sn)
 &=G(3n)+N(3n)+4
   =3G(n)+2N(n)+4,\\
 3G(Sn)+N(Sn)
 &=3G(n)+3N(n)+3+Z(N(n)).
\end{align*}
The expressions agree in both Boolean cases for $N(n)$.
The base is $0=0$.  In particular,
\eqref{eq:hybrid-ramp-triple} was proved internally; it was not
inferred from an intended floor-function interpretation of $G$.

Set
\[
 H(u)=C(N(u),G(u),u).
\]
The crucial choice is to retain $u$ when its parity is $1$.
The preceding equations give
\begin{equation}
 H(2u)=3u,\qquad H(1)=1,\qquad H(3u)=3H(u).
 \label{eq:hybrid-factor-map}
\end{equation}
To check the last identity, use $N(3u)=N(u)$ and split the Boolean
$N(u)$.  At $0$, both sides are $3G(u)$ by
\eqref{eq:hybrid-ramp-triple}; at $1$, both are $3u$.
This argument does not require that an arbitrary number has already
been represented as $2v$ or $2v+1$.

\subsection{A comparator from scalar iteration}

For every fixed update $h$, equation induction on $n$, keeping $m,a$
and all parameters fixed, proves the semigroup identity
\[
 I(m+n,a;h)=I(n,I(m,a;h);h).
\]
The base is the zero computation and the step applies $h$ to the
induction equation. This identity will also be used in the later
pairing construction.


For a fixed update $h$, equation induction on $y$, with its initializer
$a$ and all parameters of $h$ fixed, proves the rotation identity
\begin{equation}
 I(Sy,a;h)=I(y,h(a);h).
 \label{eq:hybrid-rotation}
\end{equation}
Both sides have base $h(a)$, and their successor equations apply the
same $h$ to the previous equality.  Here and below $I(y,a;h)$ is
an abbreviation for one fixed finite scalar term, not quantification
over functions.

By \eqref{eq:hybrid-factor-map}, another equation induction gives
\begin{equation}
 I(y,3a;H)=3I(y,a;H).
 \label{eq:hybrid-iterator-scaling}
\end{equation}
Its base is $3a=3a$.  In the step the fixed-parameter induction
equation rewrites the new left side as
\[
 H(3I(y,a;H))=3H(I(y,a;H)).
\]
The initializer parameter stays $a$ throughout this induction.
Equation induction also gives $I(y,1;H)=1$, since $H(1)=1$.

Define
\[
 E(x)=I(x,1;u.2u),\qquad
 L(x,y)=N\bigl(I(y,E(x);H)\bigr).
\]
Then the following entire recursion specification is provable:
\begin{equation}
 L(0,y)=1,\qquad L(Sx,0)=0,\qquad L(Sx,Sy)=L(x,y).
 \label{eq:hybrid-comparator}
\end{equation}
The first equation uses $E(0)=1$ and the fixed-point identity for
$H$.  The second uses $E(Sx)=2E(x)$ and $N(2a)=0$.
For the third,
\begin{align*}
 L(Sx,Sy)
 &=N\bigl(I(Sy,2E(x);H)\bigr)\\
 &=N\bigl(I(y,H(2E(x));H)\bigr)\\
 &=N\bigl(I(y,3E(x);H)\bigr)\\
 &=N\bigl(3I(y,E(x);H)\bigr)\\
 &=N\bigl(I(y,E(x);H)\bigr)=L(x,y).
\end{align*}
The successive steps use rotation, $H(2a)=3a$, iterator scaling,
and $N(3a)=N(a)$.  Thus $L$ has the comparison equations at all
arguments internally.  Its standard-model meaning as the indicator
of $x\le y$ is motivation, not a premise of the following argument.

\subsection{Saturated iteration gives predecessor}

We now need only the three equations
\eqref{eq:hybrid-comparator}.  Define
\[
 B_m(u)=C(L(Su,m),u,Su),\qquad
 W(t,m)=I(t,0;u.B_m(u)).
\]
The subscript in $B_m$ records a fixed parameter; $B_m$ is an open
term abbreviation, not a function variable.

Numerical cases on $e$ prove the selector equation
\[
 C(e,Sa,Sb)=S C(e,a,b).
\]
Equivalently one may use equation induction on $e$, whose induction
hypothesis is unused.  Therefore
\begin{align}
 B_{Sm}(Su)
 &=C(L(SSu,Sm),Su,SSu)\notag\\
 &=C(L(Su,m),Su,SSu)=S B_m(u),
 \label{eq:hybrid-conjugacy}\\
 B_{Sm}(0)&=C(L(S0,Sm),0,S0)=C(L(0,m),0,S0)=S0.
 \label{eq:hybrid-conjugacy-base}
\end{align}

Literal-equation induction on $t$, with $m$ fixed, proves
\begin{equation}
 W(St,Sm)=S W(t,m).
 \label{eq:hybrid-saturated-shift}
\end{equation}
The base is \eqref{eq:hybrid-conjugacy-base}.  The successor step is
\begin{align*}
 W(SSt,Sm)
 &=B_{Sm}(W(St,Sm))\\
 &=B_{Sm}(S W(t,m))\\
 &=S B_m(W(t,m))=S W(St,m).
\end{align*}
Only the second line uses the current induction hypothesis, at the
same fixed $m$.

Two subsequent equation inductions give
\begin{equation}
 L(Sn,n)=0,\qquad W(n,n)=n.
 \label{eq:hybrid-diagonals}
\end{equation}
The first uses the comparator recurrence.  The second uses
\eqref{eq:hybrid-saturated-shift} in its successor step.  The
substitution $t=m=n$ is made into that already proved universally
parameterized equation; no parameter of an active induction hypothesis
is changed.  Consequently
\begin{equation}
 W(Sn,n)=B_n(n)=C(L(Sn,n),n,Sn)=n.
 \label{eq:hybrid-extra-step}
\end{equation}

Finally define the finite scalar term
\begin{equation}
 P(m)=I\bigl(m,0;u.C(L(SSu,m),u,Su)\bigr).
 \label{eq:hybrid-predecessor-term}
\end{equation}
Its zero equation is immediate.  To justify its successor equation
without assuming a separate extensionality rule for iterator bodies,
put
\[
 V(t,m)=I\bigl(t,0;u.C(L(SSu,Sm),u,Su)\bigr).
\]
The comparator recurrence identifies this update pointwise with
$B_m(u)$.  Equation induction on $t$, with $m$ fixed, gives
$V(t,m)=W(t,m)$: their base values are $0$, and the successor step
uses the pointwise update equation and then the induction equation.
Now \eqref{eq:hybrid-extra-step} yields
\[
 \boxed{\qquad P(0)=0,\qquad P(Sm)=V(Sm,m)=W(Sm,m)=m.\qquad}
\]

\subsection{Dependency boundary}

All terms in this section are built from $0,S,I$ by finite substitution.
All inductions have a single equation as their invariant.  Numerical
cases and equality transport are the already reconstructed logical
resources.  Successor injectivity may now be obtained by applying $P$
to an assumed equation $Sx=Sy$.  Equality tests, reflection of
quantifier-free formulas, pairing, and the recovery of general
primitive recursion can then be developed in the established order.
They were not used in order to manufacture predecessor.

This argument is a written internal derivation in the stated theory.
The finite implementation checks do not constitute a proof-assistant
formalization of this entire metatheorem.


% Standalone section.  Requires only amsmath and amssymb in the parent document.
% No theorem environments, macros, or bibliography keys are assumed.
% All arithmetic claims in this section are proved here.

\section{From literal-equation induction to full intuitionistic PRA}
\label{sec:eic-completion}\label{app:arithmetic}

This section supplies an unconditional completion theorem. It is adapted
with attribution from the current forward-left draft~\cite{forward}. In the
scalar application used here, every function occurring in the proof is
already a scalar term: predecessor is supplied by
Appendix~\ref{sec:hybrid-bootstrap}, the selector is
$C(n,a,b)=I(n,a;u.b)$, addition is $x+y=I(y,x;u.Su)$, and
truncated subtraction is $d(x,y)=I(y,x;u.Pu)$. All other functions below
are finite compositions of these. Thus this proof does not presuppose
that arbitrary primitive recursion has already been compiled. The QFI
conclusion first applies to the scalar language; the general PR compiler
comes afterward.  Literal-equation
induction, together with the primitive-recursive defining equations,
predecessor coverage, and separation of $0$ and $1$, proves induction for every
quantifier-free formula.  In particular, the proof does not assume a faithful
numerical equality test, stability or decidability of equality, successor
uniqueness, or any commutation law for an iterator.  We construct the equality
test and prove its reflection property for arbitrary elements of the
numerical domain.

\subsection{The exact hypotheses}

Work in ordinary intuitionistic first-order logic with equality and genuine
falsity $\bot$.  The language has one numerical sort and a symbol for each
fixed, finite primitive-recursive description.  Distinct descriptions are
not identified merely because they compute the same standard function.
Let $\mathsf{PRE}$ contain the defining equations for zero, successor,
projections, composition, and primitive recursion.  In particular, choose
descriptions $P$ and $C$ with equations
\begin{align}
 P(0)&=0, & P(Sx)&=x, \label{eq:eic-predecessor}\\
 C(0,a,b)&=a, & C(Sx,a,b)&=b. \label{eq:eic-selector}
\end{align}
Write $1=S0$.  In addition to $\mathsf{PRE}$, assume
\begin{align}
 \mathsf{Cov}:&\qquad
       \forall x\,\bigl(x=0\lor x=S(Px)\bigr),
       \label{eq:eic-coverage}\\
 \mathsf{Sep}:&\qquad 0=1\longrightarrow\bot,
       \label{eq:eic-separation}
\end{align}
and induction for each \emph{literal equation} $E(n,\vec p)$:
\begin{equation}
 \forall\vec p\,
 \left[
   \left(E(0,\vec p)\land
       \forall n\bigl(E(n,\vec p)\longrightarrow E(Sn,\vec p)\bigr)
   \right)
   \longrightarrow\forall n\,E(n,\vec p)
 \right].
 \label{eq:eic-induction}
\end{equation}
Here the schematic expression $E$ is a single equality of numerical
terms; the displayed universal quantifiers range only over numbers.
Every description and every instance of the schema is finite external
syntax.

The axiom formulation \eqref{eq:eic-induction} and the usual contextual
induction rule are interchangeable.  For example, if
\[
 \Gamma\vdash E(0,\vec p),
 \qquad
 \Gamma,E(z,\vec p)\vdash E(Sz,\vec p),
\]
where $z$ is absent from $\Gamma$ and the passive parameters, discharge the
second premise's equation and universally introduce $z$.  The axiom then
gives $\Gamma\vdash\forall n E(n,\vec p)$, which can be instantiated at any
target term.  Conversely, assume the antecedent of the axiom and apply the
contextual rule with fresh variables.  This uses the ordinary intuitionistic
logical rules and permits arbitrary, including quantified, formulas in
$\Gamma$.

\subsection{Constructor separation and two elementary tests}

Define the sign description and the countdown description by
\begin{equation}
 \operatorname{sg}(u)=C(u,0,1),
 \qquad
 K(u)=C(u,1,Pu).
 \label{eq:eic-sign-countdown}
\end{equation}
Their defining equations imply
\begin{align}
 \operatorname{sg}(0)&=0,
 &\operatorname{sg}(Sx)&=1,\label{eq:eic-sign-computation}\\
 K(0)&=1,
 &K(Sx)&=x.\label{eq:eic-countdown-computation}
\end{align}

\paragraph{Constructor laws.}
Successor is injective: applying $P$ to $Sx=Sy$ gives $x=y$ by
\eqref{eq:eic-predecessor}.  Successor never equals zero.  Indeed,
$Sx=0$ gives $1=0$ by applying $\operatorname{sg}$ and using
\eqref{eq:eic-sign-computation}; symmetry and
\eqref{eq:eic-separation} give $\bot$.  Consequently $1\ne0$, where
$A\ne B$ abbreviates $A=B\to\bot$.  None of these arguments uses
induction.

\paragraph{The sign test.}
For every $u$,
\begin{align}
 &\operatorname{sg}(u)=0\lor\operatorname{sg}(u)=1,
       \label{eq:eic-sign-range}\\
 &\operatorname{sg}(u)=0\quad\longleftrightarrow\quad u=0,
       \label{eq:eic-sign-zero}\\
 &\operatorname{sg}(u)=1\quad\longleftrightarrow\quad u\ne0.
       \label{eq:eic-sign-one}
\end{align}
To prove the range assertion, split \eqref{eq:eic-coverage}.  The branch
$u=0$ gives sign $0$, and the branch $u=S(Pu)$ gives sign $1$.
The implication $u=0\to\operatorname{sg}(u)=0$ is computation.
For its converse, the successor branch of coverage would make
$\operatorname{sg}(u)=0$ imply $1=0$, a contradiction; in that branch
explosion gives $u=0$ as well.  Similarly, if
$\operatorname{sg}(u)=1$ and $u=0$, computation gives $0=1$.
Conversely, $u\ne0$ excludes the zero branch of coverage, and the
successor branch computes the sign as $1$.

Coverage and \eqref{eq:eic-selector} also give, for all $u,v$,
\begin{equation}
 C(u,v,v)=v.
 \label{eq:eic-selector-idempotence}
\end{equation}
Both coverage branches compute to $v$, so this identity again needs no
induction.

\paragraph{The countdown test and the notation $W$.}
For every $u$,
\begin{equation}
 K(u)=0\quad\longleftrightarrow\quad u=1.
 \label{eq:eic-countdown-zero}
\end{equation}
If $u=1=S0$, \eqref{eq:eic-countdown-computation} gives $K(u)=0$.
Conversely, assume $K(u)=0$ and split coverage.  If $u=0$, then
$K(u)=1$, contradicting separation, and explosion gives $u=1$.
If $u=S(Pu)$, computation gives $K(u)=Pu$; hence $Pu=0$ and
$u=S0=1$.

If a presentation uses a description named $W$ with defining equations
\[
 W(0)=1,\qquad W(Sx)=x,
\]
then it proves $K(u)=W(u)$ for every $u$ already from coverage.
In the zero branch both sides are $1$; in the branch $u=S(Pu)$ both
sides are $Pu$.  Thus $K$ and $W$ are provably equal on the whole domain,
even if they are distinct descriptions.  This equality does not appeal to
standard-model agreement or to uniqueness of an arbitrary recurrence.

\subsection{Four elementary uses of equation induction}

Choose descriptions for addition and truncated subtraction, with recursion
on the second argument:
\begin{align}
 x+0&=x, &x+Sy&=S(x+y),\label{eq:eic-addition}\\
 d(x,0)&=x, &d(x,Sy)&=P(d(x,y)).\label{eq:eic-subtraction}
\end{align}
Using these equations and \eqref{eq:eic-induction}, we now derive
\begin{align}
 0+y&=y,\label{eq:eic-left-zero}\\
 (Sx)+y&=S(x+y),\label{eq:eic-left-successor}\\
 d(Sx,Sy)&=d(x,y),\label{eq:eic-cancel-successors}\\
 d(Sa,a)&=1.\label{eq:eic-diagonal}
\end{align}
Each induction below has a single equation as its induction formula.

For \eqref{eq:eic-left-zero}, induct on $y$.  The base is $0+0=0$.
The step is $0+Sy=S(0+y)=Sy$.

For \eqref{eq:eic-left-successor}, hold $x$ fixed and induct on $y$.
At zero both sides are $Sx$.  At a successor the induction hypothesis
gives
\[
 (Sx)+Sy=S\bigl((Sx)+y\bigr)
       =S\bigl(S(x+y)\bigr)=S(x+Sy).
\]
Combining \eqref{eq:eic-left-successor} with the defining equation for
addition gives the useful rearrangement
\begin{equation}
 (Sy)+r=y+Sr.
 \label{eq:eic-shift-successor}
\end{equation}
Indeed, both sides are $S(y+r)$.

For \eqref{eq:eic-cancel-successors}, hold $x$ fixed and induct on $y$.
The base computes as
\[
 d(Sx,S0)=P(d(Sx,0))=P(Sx)=x=d(x,0).
\]
Applying $P$ to the induction equation gives the step:
\[
 d(Sx,S(Sy))=P(d(Sx,Sy))=P(d(x,y))=d(x,Sy).
\]

Finally, induct on $a$ for \eqref{eq:eic-diagonal}.  The base is
$d(S0,0)=S0=1$.  At a successor,
\[
 d(S(Sa),Sa)=d(Sa,a)=1,
\]
using \eqref{eq:eic-cancel-successors} and the induction hypothesis.

\subsection{An equation invariant that reflects actual equality}

The remaining arithmetic ingredient is the following equation, valid for
all $x,y$:
\begin{equation}
 C\bigl(\operatorname{sg}(d(x,y)),\ x,\ y+d(x,y)\bigr)=x.
 \label{eq:eic-residual-invariant}
\end{equation}
Its meaning is that a nonzero subtraction residual reconstructs the
original argument by addition.  The selector expresses this assertion as
one literal equation, so equation induction applies directly.

Hold $x$ fixed and induct on $y$.  At $y=0$, subtraction gives residual
$x$, and \eqref{eq:eic-left-zero} gives $0+x=x$.  The desired equation
is consequently $C(\operatorname{sg}(x),x,x)=x$, which is
\eqref{eq:eic-selector-idempotence}.

For the step, assume \eqref{eq:eic-residual-invariant} at $y$ and
abbreviate
\[
 r=d(x,y),\qquad r'=Pr=d(x,Sy).
\]
These are local term abbreviations, not additional function symbols.
Split coverage for $r'$.

If $r'=0$, then $\operatorname{sg}(r')=0$, and the successor instance
of \eqref{eq:eic-residual-invariant} reduces immediately to $x=x$.

In the other branch, $r'=S(Pr')$, so constructor separation gives
$r'\ne0$.  Split coverage for $r$.  Its zero branch would imply
$r'=P0=0$, a contradiction, from which the desired step equation follows.
In its successor branch,
\[
 r=S(Pr)=Sr',\qquad \operatorname{sg}(r)=1.
\]
The induction equation therefore computes to $y+r=x$.  Meanwhile,
\eqref{eq:eic-shift-successor} gives
\[
 (Sy)+r'=y+Sr'=y+r=x.
\]
Since $r'$ is a successor, $\operatorname{sg}(r')=1$, so this is
exactly the branch selected in the successor instance of
\eqref{eq:eic-residual-invariant}.  This completes the equation induction.
All case distinctions came from coverage; no induction on an implication
or a disjunction was used.

\paragraph{Reflection at residual $1$.}
We have, for arbitrary $a,b$,
\begin{equation}
 d(Sa,b)=1\quad\longleftrightarrow\quad a=b.
 \label{eq:eic-residual-reflection}
\end{equation}
For the forward direction, substitute $x=Sa$, $y=b$ into
\eqref{eq:eic-residual-invariant}.  If $d(Sa,b)=1$, the sign is $1$
and the selector returns its second branch, giving
\[
 b+1=Sa.
\]
The defining equation for addition gives $b+1=S(b+0)=Sb$.  Thus
$Sb=Sa$, and applying $P$ gives $b=a$, hence $a=b$.
Conversely, $a=b$ and \eqref{eq:eic-diagonal} give
$d(Sa,b)=d(Sa,a)=1$.

Define the equality characteristic by the fixed description
\begin{equation}
 e(a,b)=\operatorname{sg}\bigl(K(d(Sa,b))\bigr).
 \label{eq:eic-equality-characteristic}
\end{equation}
The sign, countdown, and residual reflection lemmas give
\begin{align}
 e(a,b)=0&\quad\longleftrightarrow\quad a=b,
       \label{eq:eic-equality-reflection}\\
 e(a,b)=0&\ \lor\ e(a,b)=1.
       \label{eq:eic-equality-range}
\end{align}
In detail, $e(a,b)=0$ is equivalent to $K(d(Sa,b))=0$ by
\eqref{eq:eic-sign-zero}, which is equivalent to $d(Sa,b)=1$ by
\eqref{eq:eic-countdown-zero}, and then to $a=b$ by
\eqref{eq:eic-residual-reflection}.  The range assertion is
\eqref{eq:eic-sign-range} applied to $K(d(Sa,b))$.

The reflection property concerns the logic's actual equality throughout
the domain.  It was not assumed as part of a numerical comparison symbol,
and its proof did not restrict $a$ or $b$ to standard numerals.
It also yields decidable equality.  In the $e(a,b)=0$ branch,
\eqref{eq:eic-equality-reflection} gives $a=b$.  In the $e(a,b)=1$
branch, an assumption $a=b$ would give $e(a,b)=0$ and therefore $0=1$,
contradicting separation.  Hence
\begin{equation}
 \forall a\,\forall b\,\bigl(a=b\lor(a=b\to\bot)\bigr)
 \label{eq:eic-equality-decidable}
\end{equation}
is a theorem, not an extra hypothesis.

\subsection{Compilation of all quantifier-free formulas}

For each fixed quantifier-free formula $A(\vec p)$, construct a term
$b_A(\vec p)$ by external recursion on its finite syntax.  Take equality,
$\bot$, $\land$, $\lor$, and $\to$ as the primitive formula constructors;
negation is implication into $\bot$.  Set
\begin{align}
 b_{s=t}&=e(s,t), & b_{\bot}&=1,\label{eq:eic-boolean-atoms}\\
 b_{A\land D}&=C(b_A,b_D,1),\label{eq:eic-boolean-and}\\
 b_{A\lor D}&=C(b_A,0,b_D),\label{eq:eic-boolean-or}\\
 b_{A\to D}&=C(b_A,b_D,0).\label{eq:eic-boolean-implies}
\end{align}
Every term here is a finite composition of fixed primitive-recursive
descriptions; if desired, projections and composition give it a single
description symbol.  No internal quantification over formulas or programs
is involved.

An external induction on the syntax of $A$ proves, within the arithmetic
theory,
\begin{equation}
 \bigl(b_A=0\lor b_A=1\bigr)
 \quad\text{and}\quad
 \bigl(b_A=0\longleftrightarrow A\bigr).
 \label{eq:eic-boolean-correctness}
\end{equation}
Here are the full connective checks.  In every case $b_A=1$ implies
$\neg A$: if $A$ held, the already established implication
$A\to b_A=0$ would contradict separation.

\paragraph{Atomic formulas and falsity.}
For an equation, both claims are
\eqref{eq:eic-equality-reflection}--\eqref{eq:eic-equality-range}.
For falsity the range assertion is $1=0\lor1=1$, with its right disjunct
proved by reflexivity.  Also $1=0\to\bot$ follows from separation and
symmetry, while $\bot\to1=0$ follows by explosion.  Thus $1=0$ is
equivalent to $\bot$, as required.

\paragraph{Conjunction.}
Split $b_A=0\lor b_A=1$.  In the zero branch, $A$ holds and
$b_{A\land D}=b_D$.  Its range follows from the range of $b_D$,
and its being zero is equivalent to $D$, hence to $A\land D$ under
the known $A$.  In the one branch, $\neg A$ holds and
$b_{A\land D}=1$.  Its range is immediate.  If it were zero,
separation would give $\bot$, hence $A\land D$; conversely an
assumption $A\land D$ contradicts $\neg A$, and explosion gives
$b_{A\land D}=0$.

\paragraph{Disjunction.}
In the branch $b_A=0$, the formula $A$ holds and
$b_{A\lor D}=0$.  The range is immediate, and both $b_{A\lor D}=0$
and $A\lor D$ hold, proving their equivalence.  In the branch
$b_A=1$, we have $\neg A$ and $b_{A\lor D}=b_D$.  Its range follows
from the range of $b_D$.  If its value is zero, then $D$ and hence
$A\lor D$.  Conversely, from $A\lor D$, the $A$ alternative
contradicts $\neg A$ and yields $b_D=0$ by explosion, while the
$D$ alternative yields $b_D=0$ directly.  This proves the required
equivalence in both alternatives.

\paragraph{Implication.}
In the branch $b_A=0$, the formula $A$ holds and
$b_{A\to D}=b_D$.  Its range follows from the range of $b_D$.
If its value is zero, then $D$, which gives $A\to D$.  Conversely,
$A\to D$ and the known $A$ give $D$, hence $b_D=0$.
In the branch $b_A=1$, we have $\neg A$ and
$b_{A\to D}=0$.  Its range is immediate.  To prove $A\to D$,
assume $A$; this contradicts $\neg A$, and explosion gives $D$.
Thus both $b_{A\to D}=0$ and $A\to D$ hold in this branch, and
their equivalence follows.  This argument uses the proved two-value
range of the finite characteristic.  It does not use excluded middle
for arbitrary formulas or change the surrounding intuitionistic logic.

This completes the external syntactic induction proving
\eqref{eq:eic-boolean-correctness} for every quantifier-free formula.

\subsection{Recovery of induction and the full theory}

Let $A(n,\vec p)$ be any fixed quantifier-free formula.  Suppose an
arbitrary retained context $\Gamma$ proves its base and its induction
step:
\[
 \Gamma\vdash A(0,\vec p),
 \qquad
 \Gamma,A(z,\vec p)\vdash A(Sz,\vec p),
\]
with the usual freshness condition on $z$.  By
\eqref{eq:eic-boolean-correctness}, the base gives
$\Gamma\vdash b_A(0,\vec p)=0$.  Under the additional assumption
$b_A(z,\vec p)=0$, the same equivalence gives $A(z,\vec p)$,
the given step gives $A(Sz,\vec p)$, and the equivalence gives
$b_A(Sz,\vec p)=0$.  Apply literal-equation induction to
$b_A(n,\vec p)=0$ and then apply its reflection implication.
The result is $\Gamma\vdash A(t,\vec p)$ for every target term $t$.
Thus the theory proves the full contextual quantifier-free induction
rule, including in arbitrary quantified contexts.

\paragraph{Completion theorem.}
In the common language of fixed primitive-recursive descriptions, the
theory
\[
 \mathsf{PRE}+\mathsf{Cov}+\mathsf{Sep}+\mathsf{EI}
\]
has exactly the same sequents as ordinary first-order intuitionistic
primitive recursive arithmetic with all quantifier-free induction
instances.  The forward simulation follows from the recovered
quantifier-free induction, successor injectivity, and successor separation;
all primitive-recursive defining equations and all logical rules are
shared.  For the reverse simulation, literal-equation induction is an
instance of quantifier-free induction, and successor separation gives
$\mathsf{Sep}$.  Quantifier-free induction proves $\mathsf{Cov}$:
at zero choose its left disjunct, and at a successor use
$Sx=S(P(Sx))$, which follows from $P(Sx)=x$, for its right disjunct.

Both interpretations therefore use the full domain, defined by $x=x$,
and preserve every function symbol, equality, logical connective, and
quantifier.  Both composites have the identity isomorphism, with graph
$x=y$.  The theories are consequently identity-domain bi-interpretable
and equiconsistent.  This completion theorem is unconditional once the
literal-equation induction rule has been obtained; it has no dependency
on a classical successor-uniqueness reduction.


\section{Pairing, primitive recursion, and internal numerical composites}
\label{sec:iterator-coding}
This appendix is adapted with attribution from the progress-closure
arithmetic development~\cite{progress}. Its prerequisites have now been
proved: scalar predecessor, arithmetic reflection, and full quantifier-free
induction. Pairing is a derived proof tool; it is not part of any checker.

\subsection{Pairing from a bounded square-root search}

Once QFI is available, a familiar square-based injection gives a
particularly direct compiler. Write $a\le b$ for $d(a,b)=0$.
The usual linear-order and semiring identities follow by fixed
quantifier-free inductions from the successor, addition, and monus
equations. Multiplication itself is the scalar term
\[
                        a\cdot b=I(b,0;u.u+a).
\]
Only primitive-recursive arithmetic is involved in the following
bounded root search; there is no minimization over an unbounded domain.

For a fixed parameter $n$, put
\[
 U_n(v)=C(d((Sv)^2,n),Sv,v),\qquad
 R_n(i)=I(i,0;v.U_n(v)),\qquad q(n)=R_n(n).
\]
The update increments $v$ exactly when $(v+1)^2\le n$, and
otherwise retains it. QFI in $i$, with $n$ fixed, proves
\[
 R_n(i)\le i,\qquad R_n(i)^2\le n,\qquad
 R_n(i)=i\ \vee\ n<(S(R_n(i)))^2.              \tag{Root invariant}
\]
The base has $R_n(0)=0$. If the update increments, its test
contradicts the strict-upper-bound alternative in the third invariant.
Thus the old state equals $i$, and the new state equals $Si$;
its square bound is the increment test. If the update does not increment, the
strict upper square bound holds and the state remains fixed.
An already established strict upper square bound forces every
subsequent update to remain fixed. These are exactly the cases
needed for the three displayed invariants.

At $i=n$, either the strict upper square bound already holds or
$R_n(n)=n$, in which case $n<(n+1)^2$ supplies it. Thus
\[
                        q(n)^2\le n<(q(n)+1)^2.           \tag{Root}
\]
Monotonicity of squares shows uniqueness of a number satisfying
Root. All formulas just used are quantifier free in the scalar
language, with fixed $n$ during the induction.

Now define
\[
 \pi(a,b)=(a+b)^2+a,\qquad
 L(n)=d(n,q(n)^2),\qquad H(n)=d(q(n),L(n)).
\]
For $k=a+b$, one has $k^2\le k^2+a<(k+1)^2$ since $a\le k$.
Root uniqueness gives $q(\pi(a,b))=a+b$, and therefore
\[
                    L(\pi(a,b))=a,\qquad H(\pi(a,b))=b.
\]
Surjectivity of $\pi$ is unnecessary. The original numerical domain
is never replaced by the range of this coding; pairing is used only
to store the finite state of a fixed primitive recursion. This
standard square pairing is a convenient reconstruction here, not
a claimed new arithmetic basis.

\subsection{Compiling each primitive recursion}

Suppose the fixed target description is
\[
 f(0,\vec p)=g(\vec p),\qquad
 f(Sn,\vec p)=h(n,f(n,\vec p),\vec p).
\]
Assume by external structural induction that the earlier $g,h$
already have scalar translations $g^*,h^*$. Define
\[
 W(u,\vec p)=\pi(SL(u),h^*(L(u),H(u),\vec p)),
\]
\[
 V(n,\vec p)=I(n,\pi(0,g^*(\vec p));u.W(u,\vec p)),
 \qquad f^*(n,\vec p)=H(V(n,\vec p)).
\]
The update has only its state variable bound, with every $p_i$
fixed. Equation induction proves $L(V(n,\vec p))=n$.
The two projection equations and scalar computation then prove
both defining equations for $f^*$ internally. Initial functions
and composition translate directly. This completes the translation
of every fixed finite PR description without adding a universal
evaluator.

\subsection{Why both numerical round trips are internal}

Two fixed scalar iterators with the same seed and provably equal
updates are internally equal at all counters: use literal-equation
induction comparing their values. Their base agrees, and their
successor equations plus congruence prove the step. The same
uniqueness argument works for ordinary primitive recursion.

For an ordinary PR description, induct externally on its finite
construction. The double translation satisfies its original
defining equations, so internal equation induction identifies it
with the original operation. For a native scalar description,
earlier translated updates agree by the external induction
hypothesis; scalar uniqueness identifies the double-translated
iterator. These proofs hold at arbitrary internal counters.

Thus the domain relation $x=y$ is provably a total bijection
preserving every fixed numerical operation in both composites.
The formula constructions and round trips in
Section~\ref{sec:choice-introduction} give the corresponding formula
composites. Both interpretations preserve falsity, taking native
$0=1$ to an ordinary contradiction and ordinary $\bot$ to native
$0=1$. Finite proof translations therefore establish equiconsistency.


\clearpage
\section{The complete selected proof checker}
\label{sec:checker-listing}
The listing includes the whole standalone program; it imports no project code.
\begin{lstlisting}[language=Python]
"""Standalone five-pattern clausal iPRA checker.

Terms are zero, bound positive indices, fresh negative levels, [0,t]
(successor), and [n,a,h] (iterator). Zero iterators are normalized away.
Formulas are [s,t] and [B] (paired binder). Contexts contain clauses.
A validated reflexive goal is a leaf regardless of its unused payload.
"""
import json
import sys


def require(condition):
    if not condition:
        raise ValueError('invalid certificate')


def term(a, scope, depth=0):
    if type(a) is int:
        return -scope <= a <= depth
    return (type(a) is list and len(a) in (2, 3)
            and (a[0] != 0 if len(a) == 3 else a[0] == 0)
            and all(term(t, scope, depth + (len(a) == 3 and i == 2))
                    for i, t in enumerate(a)))


def formula(a, scope, depth=0):
    return (type(a) is list
            and ((len(a) == 2 and all(term(t, scope, depth) for t in a))
                 or (len(a) == 1 and formula(a[0], scope, depth + 2))))


def opening(a, value, depth=0):
    if type(a) is int:
        return value if a == depth + 1 else a - (a > depth + 1)
    size = len(a)
    result = [opening(t, value, depth + (2 if size == 1 else int(size == 3 and i == 2)))
              for i, t in enumerate(a)]
    return result[1] if size == 3 and result[0] == 0 else result


def prove(a, p, g, n):
    require(formula(a, n))
    if len(a) == 2 and a[0] == a[1]:
        return
    match p:
        case ['v', c, b, h, t, major, minor]:
            require(len(a) == 2 and a[1] == opening(h, t))
            prove([b if c == 0 else [c, b, h], t], major, g, n)
            prove([[[0, c], b, h], a[0]], minor, g, n)
        case ['u', index, t, premise]:
            require(type(index) is int and 0 <= index < len(g) and term(t, n))
            antecedent, body = g[index]
            require(a == opening(body, t))
            prove(antecedent, premise, g, n)
        case ['e', body, major, branch]:
            prove([body], major, g, n)
            body = opening(body, -n - 1, 1)
            prove(a, branch, g + [(opening(body, 0), body)], n + 1)
        case ['i', t, branch]:
            require(len(a) == 1 and term(t, n))
            body = opening(a[0], t, 1)
            hypothesis = opening(body, -n - 1 if len(a[0]) == 2 and t == 0 else 0)
            prove(opening(body, [0, -n - 1]), branch, g + [([0, 0], hypothesis)], n + 1)
        case _:
            raise ValueError('invalid proof')


def check(c):
    require(type(c) is dict and set(c) == {'goal', 'proof'})
    prove(c['goal'], c['proof'], [], 0)


if __name__ == '__main__':
    try:
        check(json.load(sys.stdin))
    except Exception:
        print('REJECT')
        sys.exit(1)
    print('ACCEPT')
\end{lstlisting}

\clearpage
\begin{thebibliography}{12}
\small\raggedright
\bibitem{baseline} V. Reshetnikov, \emph{ProveIt: Concise presentations of
intuitionistic primitive recursive arithmetic}, snapshot
\file{1eb151b8f40800aec569361e827d934cb77bd3e1}, accessed November 11, 2026.
\url{https://github.com/VladimirReshetnikov/ProveIt/tree/1eb151b8f40800aec569361e827d934cb77bd3e1/Logic/PrimitiveRecursiveArithmetic/ConciseIPRA}.
\bibitem{iteration} ProveIt incoming package,
\emph{Equality from iteration}, \file{ipra_equality_from_iteration_2026-11-11.zip},
in the snapshot of reference~\cite{baseline}. The zero-normal single
evaluation rule and identity-frame argument are inherited from this package.
\bibitem{rawclausal} ProveIt incoming package,
\emph{Five rules for iPRA with ordinary substitution},
\file{ipra_raw_clausal_2026-11-11.zip}, in the snapshot of
reference~\cite{baseline}. The raw CE5 control and clause translation
are inherited from this package.
\bibitem{clausal} ProveIt, \file{Research/clausal_canonical/}, in the snapshot
of reference~\cite{baseline}. Canonical clause presentation and its complete checker.
\bibitem{normal} ProveIt, \file{Lean/ConciseIPRA/ZeroNormalization.lean}
and \file{NativeZeroNormal.lean}, in the snapshot of reference~\cite{baseline}.
The endpoint \file{nzprf_iff_nprf} compares the native kernel with its
zero-normal presentation.
\bibitem{bridge} ProveIt, \file{Lean/ConciseIPRA/BiInterpretation.lean}
and \file{NativeBiInterpretation.lean}, in the snapshot of
reference~\cite{baseline}. Full native/$T_I$/iPRA interpretation chain.
\bibitem{choice} ProveIt, \file{Research/certified_zero/reference/},
in the snapshot of reference~\cite{baseline}. Inherited choice basis and
ordinary arithmetic bridge, reproduced here from the attributed copy
in the raw-clausal package.
\bibitem{erasure} ProveIt contributors, \emph{Internal scalar predecessor
bootstrap}, \file{Research/certified_zero/reference/scalar_predecessor.tex},
in the snapshot of reference~\cite{baseline}; reproduced in this report.
\bibitem{forward} ProveIt contributors, \emph{From literal-equation
induction to full intuitionistic PRA},
\file{Research/certified_zero/reference/equality_reflection.tex},
in the snapshot of reference~\cite{baseline}; reproduced in this report.
\bibitem{progress} ProveIt contributors, \emph{Pairing, primitive recursion,
and internal numerical composites},
\file{Research/certified_zero/reference/primitive_recursion.tex},
in the snapshot of reference~\cite{baseline}; reproduced in this report.
\bibitem{canonical} ProveIt, \file{Research/canonical_recursors/},
in the snapshot of reference~\cite{baseline}. Canonical syntax and the
516-byte compact checker reported for the separate codec.
\bibitem{computational} ProveIt, \file{Research/computational_equality/},
in the snapshot of reference~\cite{baseline}. Refl/Eval equality interface;
formal comparison in \file{Lean/ConciseIPRA/NativeEval.lean}.
\end{thebibliography}
\end{document}
